% =====================================================================
% 主文档：央行书面沟通影响收益还是波动？
% ---------------------------------------------------------------------
% 编译：xelatex main.tex  （需两次，以生成交叉引用）
% 依赖：ctex、booktabs、threeparttable、amsmath、geometry、chngcntr、
%       hyperref
% 结构：abstract.tex 为摘要（含英文），ch1--ch6 为正文，
%       references.tex 为参考文献。
% =====================================================================

\documentclass[12pt, a4paper]{ctexart}

\usepackage[margin=2.5cm]{geometry}
\usepackage{amsmath}
\usepackage{booktabs}
\usepackage{threeparttable}
\usepackage{chngcntr}
\usepackage[colorlinks=true, linkcolor=black, citecolor=black, urlcolor=blue]{hyperref}

% 让表、图、公式按节编号（正文引用为 表 5.2、式 (4.3) 这类格式）
\counterwithin{table}{section}
\counterwithin{figure}{section}
\counterwithin{equation}{section}

\linespread{1.35}
\setlength{\parskip}{0.25em}

\title{\textbf{央行书面沟通影响收益还是波动？}\\[4pt]
\large 来自 102 期《货币政策执行报告》文本情绪的证据}
\author{}
\date{}

\begin{document}

\maketitle

% =====================================================================
% 摘要与关键词
% ---------------------------------------------------------------------
% 用法：放入主文档 \begin{document} 之后、正文之前：
%     % =====================================================================
% 摘要与关键词
% ---------------------------------------------------------------------
% 用法：放入主文档 \begin{document} 之后、正文之前：
%     % =====================================================================
% 摘要与关键词
% ---------------------------------------------------------------------
% 用法：放入主文档 \begin{document} 之后、正文之前：
%     \input{abstract}
% 适用 ctexart / ctexbook。若需英文摘要，把下面的 english 部分放在正文末尾。
% 依赖：无额外宏包（ctex 自带 abstract 环境）
% =====================================================================

\begin{abstract}
\noindent
央行《货币政策执行报告》是市场识别政策意图的主要书面渠道。已有研究确认报告文本
情绪对股市整体收益具有解释力，并指出真正起作用的是面向未来的政策指引而非宏观回顾；
但这些结论多以宽基指数为对象，较少讨论不同制度阶段是否一致。

本文以 2001 年第一季度至 2026 年第二季度共 102 期报告为样本，采用中文金融情感词典
与情感单元法，分别度量政策指引段与宏观回顾段的情绪；以 16 个申万一级行业的等权组合
构造事件窗收益与波动指标，考察报告发布后一日与五日内的市场反应。

估计结果显示，政策指引段情绪对发布后的收益方向没有稳健影响，却显著降低事件日的市场
波动：情绪每提高一个标准差，日内振幅下降 15 至 28 个基点（取决于是否包含金融危机年份），
控制发布前波动状态后为 15 至 19 个基点；该结果在 28 行业口径、500 次置换检验与
政策操作控制下均保持稳健，而宏观回顾段情绪无此效应。分时段来看，波动效应集中于
2014 年以前；2014 年以后情绪与发布后收益转为负相关，但这一方向性结果对危机年份高度
敏感，本文仅将其作为提示性证据。

上述发现意味着，央行书面沟通的作用首先体现在收敛市场分歧，而非改变价格方向，且这种
作用随货币政策框架的调整而变化。若把不同时期、不同文本段落视为同质样本处理，方向相反
的两类效应会相互抵消，容易得出沟通无效的错误判断。
\end{abstract}

\vspace{0.5em}
\noindent\textbf{关键词：}货币政策执行报告；文本情绪；市场波动；结构断点；事件研究

\vspace{0.3em}
\noindent\textbf{JEL 分类号：}E52；E58；G14

% ---------------------------------------------------------------------
% 英文摘要（按需保留；若期刊或赛事不要求可删去）
% ---------------------------------------------------------------------
\clearpage
\renewcommand{\abstractname}{Abstract}
\begin{abstract}
\noindent
The Monetary Policy Report (MPR) published by the People's Bank of China is the
primary written channel through which policy intentions are communicated to the
market. Existing studies have established that the tone of the MPR explains
aggregate stock market returns, and that the forward-looking policy guidance
rather than the macroeconomic review drives this reaction. Most of this evidence,
however, rests on broad market indices, and whether it holds across institutional
periods has received little attention. Using 102 quarterly reports from 2001Q1 to
2026Q2 and a Chinese financial sentiment dictionary, we measure the tone of the two
segments separately. Market reactions are measured with event-window returns and
volatility constructed from an equal-weighted portfolio of 16 Shenwan level-1
industries.

We find no robust effect of policy guidance tone on the direction of post-release
returns, but a significant reduction in announcement-day volatility: a one
standard deviation increase in tone lowers the intraday range by 15 to 28 basis
points, depending on whether crisis years are included, and by 15 to 19 basis
points after controlling for pre-announcement volatility. The result survives an
alternative 28-industry market definition, a 500-draw permutation test, and
controls for reserve requirement ratio cuts, LPR cuts and reverse repo rate cuts,
while the tone of the macroeconomic review shows no such effect. The volatility
effect is concentrated before 2014; after 2014 the relation between tone and
returns turns negative, but this directional result is highly sensitive to crisis
years and is therefore reported as suggestive evidence.

These findings suggest that written central bank communication primarily reduces
disagreement rather than moving prices, and that its effect is not constant but
varies with the monetary policy framework. Treating heterogeneous periods and text
segments as a homogeneous sample may cancel out effects of opposite sign and lead
to the misleading conclusion that communication does not matter.
\end{abstract}

\vspace{0.5em}
\noindent\textbf{Keywords:} Monetary Policy Report; textual sentiment; market volatility;
structural break; event study

% ---------------------------------------------------------------------
% 写作依据（仅供团队核对，不进入排版）
%   "面向未来的政策指引而非宏观回顾"        → 姜富伟等（2021）的核心结论
%   前瞻性/回顾性分段的处理思路              → 王辉等（2023）；蒋海等（2025）
%   文本度量对词典与设定的敏感性（稳健性）    → 林建浩、孙乐轩（2025）
%   "信息效应"而非单纯冲击的分析框架          → 王熙、郑梦圆（2021）
%   行业事件窗与行业面板做法                  → 游庭玮（2026）；王春芝（2022）；周亮（2018）
% ---------------------------------------------------------------------

% 适用 ctexart / ctexbook。若需英文摘要，把下面的 english 部分放在正文末尾。
% 依赖：无额外宏包（ctex 自带 abstract 环境）
% =====================================================================

\begin{abstract}
\noindent
央行《货币政策执行报告》是市场识别政策意图的主要书面渠道。已有研究确认报告文本
情绪对股市整体收益具有解释力，并指出真正起作用的是面向未来的政策指引而非宏观回顾；
但这些结论多以宽基指数为对象，较少讨论不同制度阶段是否一致。

本文以 2001 年第一季度至 2026 年第二季度共 102 期报告为样本，采用中文金融情感词典
与情感单元法，分别度量政策指引段与宏观回顾段的情绪；以 16 个申万一级行业的等权组合
构造事件窗收益与波动指标，考察报告发布后一日与五日内的市场反应。

估计结果显示，政策指引段情绪对发布后的收益方向没有稳健影响，却显著降低事件日的市场
波动：情绪每提高一个标准差，日内振幅下降 15 至 28 个基点（取决于是否包含金融危机年份），
控制发布前波动状态后为 15 至 19 个基点；该结果在 28 行业口径、500 次置换检验与
政策操作控制下均保持稳健，而宏观回顾段情绪无此效应。分时段来看，波动效应集中于
2014 年以前；2014 年以后情绪与发布后收益转为负相关，但这一方向性结果对危机年份高度
敏感，本文仅将其作为提示性证据。

上述发现意味着，央行书面沟通的作用首先体现在收敛市场分歧，而非改变价格方向，且这种
作用随货币政策框架的调整而变化。若把不同时期、不同文本段落视为同质样本处理，方向相反
的两类效应会相互抵消，容易得出沟通无效的错误判断。
\end{abstract}

\vspace{0.5em}
\noindent\textbf{关键词：}货币政策执行报告；文本情绪；市场波动；结构断点；事件研究

\vspace{0.3em}
\noindent\textbf{JEL 分类号：}E52；E58；G14

% ---------------------------------------------------------------------
% 英文摘要（按需保留；若期刊或赛事不要求可删去）
% ---------------------------------------------------------------------
\clearpage
\renewcommand{\abstractname}{Abstract}
\begin{abstract}
\noindent
The Monetary Policy Report (MPR) published by the People's Bank of China is the
primary written channel through which policy intentions are communicated to the
market. Existing studies have established that the tone of the MPR explains
aggregate stock market returns, and that the forward-looking policy guidance
rather than the macroeconomic review drives this reaction. Most of this evidence,
however, rests on broad market indices, and whether it holds across institutional
periods has received little attention. Using 102 quarterly reports from 2001Q1 to
2026Q2 and a Chinese financial sentiment dictionary, we measure the tone of the two
segments separately. Market reactions are measured with event-window returns and
volatility constructed from an equal-weighted portfolio of 16 Shenwan level-1
industries.

We find no robust effect of policy guidance tone on the direction of post-release
returns, but a significant reduction in announcement-day volatility: a one
standard deviation increase in tone lowers the intraday range by 15 to 28 basis
points, depending on whether crisis years are included, and by 15 to 19 basis
points after controlling for pre-announcement volatility. The result survives an
alternative 28-industry market definition, a 500-draw permutation test, and
controls for reserve requirement ratio cuts, LPR cuts and reverse repo rate cuts,
while the tone of the macroeconomic review shows no such effect. The volatility
effect is concentrated before 2014; after 2014 the relation between tone and
returns turns negative, but this directional result is highly sensitive to crisis
years and is therefore reported as suggestive evidence.

These findings suggest that written central bank communication primarily reduces
disagreement rather than moving prices, and that its effect is not constant but
varies with the monetary policy framework. Treating heterogeneous periods and text
segments as a homogeneous sample may cancel out effects of opposite sign and lead
to the misleading conclusion that communication does not matter.
\end{abstract}

\vspace{0.5em}
\noindent\textbf{Keywords:} Monetary Policy Report; textual sentiment; market volatility;
structural break; event study

% ---------------------------------------------------------------------
% 写作依据（仅供团队核对，不进入排版）
%   "面向未来的政策指引而非宏观回顾"        → 姜富伟等（2021）的核心结论
%   前瞻性/回顾性分段的处理思路              → 王辉等（2023）；蒋海等（2025）
%   文本度量对词典与设定的敏感性（稳健性）    → 林建浩、孙乐轩（2025）
%   "信息效应"而非单纯冲击的分析框架          → 王熙、郑梦圆（2021）
%   行业事件窗与行业面板做法                  → 游庭玮（2026）；王春芝（2022）；周亮（2018）
% ---------------------------------------------------------------------

% 适用 ctexart / ctexbook。若需英文摘要，把下面的 english 部分放在正文末尾。
% 依赖：无额外宏包（ctex 自带 abstract 环境）
% =====================================================================

\begin{abstract}
\noindent
央行《货币政策执行报告》是市场识别政策意图的主要书面渠道。已有研究确认报告文本
情绪对股市整体收益具有解释力，并指出真正起作用的是面向未来的政策指引而非宏观回顾；
但这些结论多以宽基指数为对象，较少讨论不同制度阶段是否一致。

本文以 2001 年第一季度至 2026 年第二季度共 102 期报告为样本，采用中文金融情感词典
与情感单元法，分别度量政策指引段与宏观回顾段的情绪；以 16 个申万一级行业的等权组合
构造事件窗收益与波动指标，考察报告发布后一日与五日内的市场反应。

估计结果显示，政策指引段情绪对发布后的收益方向没有稳健影响，却显著降低事件日的市场
波动：情绪每提高一个标准差，日内振幅下降 15 至 28 个基点（取决于是否包含金融危机年份），
控制发布前波动状态后为 15 至 19 个基点；该结果在 28 行业口径、500 次置换检验与
政策操作控制下均保持稳健，而宏观回顾段情绪无此效应。分时段来看，波动效应集中于
2014 年以前；2014 年以后情绪与发布后收益转为负相关，但这一方向性结果对危机年份高度
敏感，本文仅将其作为提示性证据。

上述发现意味着，央行书面沟通的作用首先体现在收敛市场分歧，而非改变价格方向，且这种
作用随货币政策框架的调整而变化。若把不同时期、不同文本段落视为同质样本处理，方向相反
的两类效应会相互抵消，容易得出沟通无效的错误判断。
\end{abstract}

\vspace{0.5em}
\noindent\textbf{关键词：}货币政策执行报告；文本情绪；市场波动；结构断点；事件研究

\vspace{0.3em}
\noindent\textbf{JEL 分类号：}E52；E58；G14

% ---------------------------------------------------------------------
% 英文摘要（按需保留；若期刊或赛事不要求可删去）
% ---------------------------------------------------------------------
\clearpage
\renewcommand{\abstractname}{Abstract}
\begin{abstract}
\noindent
The Monetary Policy Report (MPR) published by the People's Bank of China is the
primary written channel through which policy intentions are communicated to the
market. Existing studies have established that the tone of the MPR explains
aggregate stock market returns, and that the forward-looking policy guidance
rather than the macroeconomic review drives this reaction. Most of this evidence,
however, rests on broad market indices, and whether it holds across institutional
periods has received little attention. Using 102 quarterly reports from 2001Q1 to
2026Q2 and a Chinese financial sentiment dictionary, we measure the tone of the two
segments separately. Market reactions are measured with event-window returns and
volatility constructed from an equal-weighted portfolio of 16 Shenwan level-1
industries.

We find no robust effect of policy guidance tone on the direction of post-release
returns, but a significant reduction in announcement-day volatility: a one
standard deviation increase in tone lowers the intraday range by 15 to 28 basis
points, depending on whether crisis years are included, and by 15 to 19 basis
points after controlling for pre-announcement volatility. The result survives an
alternative 28-industry market definition, a 500-draw permutation test, and
controls for reserve requirement ratio cuts, LPR cuts and reverse repo rate cuts,
while the tone of the macroeconomic review shows no such effect. The volatility
effect is concentrated before 2014; after 2014 the relation between tone and
returns turns negative, but this directional result is highly sensitive to crisis
years and is therefore reported as suggestive evidence.

These findings suggest that written central bank communication primarily reduces
disagreement rather than moving prices, and that its effect is not constant but
varies with the monetary policy framework. Treating heterogeneous periods and text
segments as a homogeneous sample may cancel out effects of opposite sign and lead
to the misleading conclusion that communication does not matter.
\end{abstract}

\vspace{0.5em}
\noindent\textbf{Keywords:} Monetary Policy Report; textual sentiment; market volatility;
structural break; event study

% ---------------------------------------------------------------------
% 写作依据（仅供团队核对，不进入排版）
%   "面向未来的政策指引而非宏观回顾"        → 姜富伟等（2021）的核心结论
%   前瞻性/回顾性分段的处理思路              → 王辉等（2023）；蒋海等（2025）
%   文本度量对词典与设定的敏感性（稳健性）    → 林建浩、孙乐轩（2025）
%   "信息效应"而非单纯冲击的分析框架          → 王熙、郑梦圆（2021）
%   行业事件窗与行业面板做法                  → 游庭玮（2026）；王春芝（2022）；周亮（2018）
% ---------------------------------------------------------------------


\newpage
\tableofcontents
\newpage

% =====================================================================
% 第 1 章  引言
% ---------------------------------------------------------------------
% 编译说明同前：ctexart 下用 \section，ctexbook/report 下改为 \chapter。
% =====================================================================

\section{引言}

中国人民银行每季度发布一份《中国货币政策执行报告》。这份文件既复盘上一季度的货币信贷
与金融市场运行，也交代下一阶段的政策取向，篇幅通常在五十页上下。对市场而言，它是少数
几份能够同时提供数据与判断的官方文本；对研究者而言，它提供了连续二十余年、结构相对
稳定的书面沟通样本。一个自然的问题是：报告里的话，市场听进去了吗？

已有的回答大多集中在价格方向上。姜富伟等（2021）以四类宽基指数为对象，发现报告文本
情绪改善伴随发布后一日的正收益，并进一步指出，真正引起反应的是反映政策指引的部分，
而非回顾宏观经济基本面的部分。此后的一系列研究把这一思路延伸到不同场景：许魁（2023）
考察分析师报告在传导中的作用，王辉等（2023）关注预期引导与银行风险，蒋海等（2025）
提取前瞻性情绪解释流动性囤积，周阳、綦岚鑫（2024）构建领域词典检验股票、债券与外汇
市场的反应。这些工作共同确立了一个判断：报告文本含有市场尚未定价的信息。

但仍有三处可以推进。第一，结论多停留在市场总量，行业维度只作为稳健性出现，缺少对
“文本影响的是哪一类资产、哪一阶矩”的区分。第二，分段方式不尽相同，有的按句子内容
分类，有的按章节切分，口径差异使结论难以直接比较。第三，样本多止于 2019 年前后，
2020 年之后的非常规货币政策阶段较少被纳入，效应是否稳定缺乏检验。

推进这些问题会碰到一个技术障碍。报告级情绪每期只有一个取值，在行业面板中加入季度
固定效应后，它与时间虚拟变量完全共线，系数无法识别。这不是样本长度的问题，而是变量
层级与模型设定的必然结果。要在有时间固定效应的框架下识别情绪，只能依赖情绪与行业
特征的交互；而要在时间序列上识别，则必须把观测限制在信息到达的附近。

本文选择后者，把识别建立在报告发布日的事件窗上。具体而言，以 2001 年第一季度至
2026 年第二季度共 102 期报告为样本，用中文金融情感词典与情感单元法分别度量政策指引段
与宏观回顾段的情绪；以 16 个连续可得的申万一级行业构造等权市场组合，计算发布后一日、
五日以及发布前后的累计收益与日内振幅；在此基础上检验情绪对收益方向与波动的影响，
并按 2014 年前后分组考察效应的稳定性。

结果可以概括为三句。政策指引段的情绪对发布后的收益方向没有稳健影响，系数接近于零；
同一情绪显著降低发布日的波动，控制事前波动后依然成立；把样本按 2014 年切开后，
波动效应集中在前期，方向效应出现在后期且为负，但后者对个别年份较为敏感，本文仅将其
作为提示性证据。

相对于既有文献，本文的边际贡献有三。其一，把被解释变量从收益扩展到波动，并用日内
振幅与绝对收益两个代理交叉验证，说明文本的作用首先体现在收敛分歧而非推动价格。
其二，按章节结构切分指引段与回顾段，在保留姜富伟等（2021）核心判断的同时，给出了
一种可复现的替代口径。其三，引入制度断点并配合断点年份敏感性、替代词典与置换检验，
对结论的稳健性作出较为完整的交代。

需要说明的是，本文没有估计“哪些行业对政策指引更敏感”。受数据条件限制，行业特征
（估值、杠杆、利率敏感度）的覆盖期偏短，无法支撑稳定的交互估计，这一扩展留待后续
研究。相应地，本文的结论限定在市场层面，行业数据的作用是构造市场组合与提供稳健性
对照。

全文余下部分安排如下。第 2 章回顾央行沟通、文本情绪度量与行业传导三支文献，并提出
待检验的假说。第 3 章说明数据来源、文本处理与变量构造。第 4 章讨论识别策略与模型
设定。第 5 章报告实证结果与稳健性检验。第 6 章讨论机制、政策含义与局限，并给出结论。

% =====================================================================
% 第 2 章  文献回顾与研究假说
% ---------------------------------------------------------------------
% 编译说明同前。
% =====================================================================

\section{文献回顾与研究假说}

与本文相关的研究可以归入三支：央行书面沟通的市场效应、文本情绪的度量方法、
货币政策的行业传导。三支文献解决的问题不同，留下的空隙也不同，本章分别梳理，
最后落到本文的定位与待检验的假说。

\subsection{央行书面沟通的市场效应}

国外文献较早把央行声明当作可度量的对象。Lucca 和 Trebbi（2009）用自动化方法解析
美联储公开市场委员会声明，考察其信息如何进入资产价格；Shapiro 和 Wilson（2022）
则从声明文本中反推政策目标的权重，把文本从“语气”推进到“偏好”。这类研究的共同
前提是，央行文字中包含着关于未来政策路径的信息，而市场会对这些信息定价。

中文文献的路径相似，但对象换成了《货币政策执行报告》。姜富伟等（2021）是其中
最直接的一篇。他们以 2001 年至 2018 年的报告为样本，用中文金融情感词典计算文本
情绪、相似度与可读性，发现情绪改善伴随发布后一日的指数正收益，相似度上升伴随波动
下降，而可读性没有显著作用。更重要的是，他们把情绪拆成“政策指引”与“宏观基本面”
两部分，发现只有前者显著，理由是报告中的宏观数据在发布前已由统计局公布，市场早已
作出反应。

沿这一思路，后续研究从不同角度补充。许魁（2023）引入分析师报告，检验其在情绪传导
中的中介作用，发现分析师情绪与报告情绪高度相关，紧缩期的传导更强。王辉等（2023）
把报告拆为前瞻性与回顾性内容，构建“市场解读偏离”指标，考察预期引导对银行系统性
风险的影响。蒋海等（2025）用机器学习方法提取前瞻性情绪，解释商业银行的流动性囤积。
王熙和郑梦圆（2021）把语调指数纳入结构向量自回归，控制信息效应后“价格之谜”消失，
从宏观侧证实了文本中的信息成分。周阳、綦岚鑫（2024）则自建适配报告的领域词典，
同时考察股票、债券与外汇市场，发现市场之间存在溢出，溢出会放大或削弱情绪效应。

这些研究给出的判断大体一致：报告文本含有信息，且信息主要来自面向未来的部分。
但结论多建立在市场总量上，行业维度往往只作为稳健性出现；样本也多在 2019 年前后
截止，2020 年之后的阶段较少被检验。

\subsection{文本情绪的度量方法}

度量方法大致有三代。第一代是词典法，把文本中的情感词按极性计数再加总。姜富伟等
（2021）使用的中文金融情感词典属于这一代，其优点是口径透明、易于复现，缺点是无法
处理语境。周阳、綦岚鑫（2024）对此作了改进，在词典中纳入否定词与程度副词，使同一
情感词在不同修饰下得到不同权重。汪禹和吴彬（2025）则从词典构建入手，用种子词加点
互信息扩充、再以形式概念分析筛选主题词，为政策文本提供了一个可复制的建库流程。

第二代是机器学习与预训练模型。纪钧尧等（2026）用 FinBERT 挖掘十二个经济体央行文本
的情绪，发现情绪对汇率波动的影响强于对汇率水平的影响；蒋海等（2025）用长短期记忆
网络标注语料，再训练分类器扩展到全样本。这一代方法在语义识别上更强，但需要标注
成本，且模型版本会影响结果。

第三代是大语言模型。廖海龙（2026）用提示词工程从报告中提取情绪因子并接入量化择时
框架；林建浩和孙乐轩（2025）系统梳理了这一路径的原理与应用，同时提醒三点局限：
输出的可复现性、对提示方式的敏感性，以及黑箱带来的可解释性问题。

本文选择词典法作为主口径，原因有二。一是与姜富伟等（2021）保持可比，便于判断差异
来自样本还是来自度量；二是词典法的误差方向较为清楚，测量误差通常导致系数向零收缩，
便于解释。为回应林建浩和孙乐轩（2025）提出的可复现性问题，本文另建一套货币政策
领域词典作稳健性，并用置换检验考察结果对随机配对的敏感性。

\subsection{货币政策的行业传导}

行业层面的研究更多集中在传统货币政策变量上。周亮（2018）用 M2、SHIBOR 与 M1-M2
划分货币周期，考察行业轮动策略的效果，发现以 M2 划分时策略表现最好。游庭玮（2026）
以美联储政策利率为解释变量，用申万一级行业的日度与周度面板估计冲击效应，行业固定
效应加行业维度聚类，并通过控制市场整体收益识别行业独立冲击，结论是高估值行业对
外部利率更敏感。王春芝（2022）用事件研究法考察货币政策公告前的股价预漂移，发现
报告发布前数个交易日已存在显著的正向异常收益。

把文本与行业结合的研究相对少。陈妍玲等（2024）用报告的情绪、相似度与可读性构建
沟通指数，考察金融业及其子行业的风险承担与风险溢出，发现前瞻性报告的文本情绪显著
提高子行业间的风险溢出。这已经是行业层面的证据，但结果变量是风险而非收益，且行业
范围限于金融业。

行业层面的一个技术困难值得单独说明。报告级情绪每期只有一个取值，若在面板中加入
季度固定效应，它与时间虚拟变量共线，系数不可识别。游庭玮（2026）的处理办法是控制
市场整体收益，以此剥离共同冲击；本文则把识别放到事件窗内，用同一时点的横截面差异
替代固定效应的作用。两种做法针对的是同一个问题：如何在不牺牲时间维度的前提下控制
共同因素。

\subsection{评述与本文的位置}

综合三支文献，可以看出一处明确的空隙。文本情绪的市场效应已被反复验证，但验证集中
在收益方向；行业维度的证据要么停留在总量，要么局限于金融业的风险指标；度量口径与
样本时期的差异又使结论难以横向比较。换言之，已有研究回答了“文本是否有信息”，
但对“信息作用于价格的哪一阶矩”“这种作用是否随时间改变”讨论不足。

本文围绕这两点展开。与姜富伟等（2021）相比，本文把样本延长至 2026 年，把被解释
变量扩展到波动，并引入制度断点；与周阳、綦岚鑫（2024）相比，本文只做股票市场，
但在分段与时期两个维度上作了更细的处理；与游庭玮（2026）相比，本文的解释变量是
文本而非政策利率，因而更接近“沟通”而非“操作”的效应。

\subsection{研究假说}

根据上述讨论，提出三个待检验的假说。

\textbf{H1（方向）}：政策指引段的文本情绪越高，报告发布后的市场收益越高。
这一假说直接对应姜富伟等（2021）的结论，用于检验本文的度量与样本能否复现既有发现。
若 H1 不成立，需要区分是度量差异、样本差异，还是效应本身不稳定。

\textbf{H2（波动）}：政策指引段的文本情绪越高，报告发布日的市场波动越低；
宏观回顾段的情绪不具有这一作用。其依据是，文本提供的是关于政策路径的信息，
信息越明确，投资者分歧越小，交易行为越趋同；而回顾段的内容在发布前已被市场知悉，
不构成新信息。

\textbf{H3（断点）}：情绪效应在 2014 年前后存在差异。2013 年的流动性紧张、2014 年
利率市场化的推进、2015 年汇率形成机制的调整，以及报告自身章节结构的改写，都可能
改变市场解读文本的方式。这一假说不预设变化的方向，只检验是否存在结构性差异。

三个假说分别对应第 5 章的三组结果：H1 对应收益方程，H2 对应波动方程与双段回归，
H3 对应交互项与断点敏感性检验。

% =====================================================================
% 第 3 章  数据与度量
% ---------------------------------------------------------------------
% 编译说明与第 5 章一致：ctexart 下用 \section，ctexbook/report 下改为 \chapter。
% 依赖宏包：booktabs、threeparttable、amsmath
% =====================================================================

\section{数据与度量}

本章说明三类数据的来源与构造：报告文本及其情绪指标、行业行情与市场组合、
事件日历与收益窗口。三者的对齐方式决定了后续识别是否成立，因此对时间戳的处理
说明得较为详细。

\subsection{报告文本与情绪度量}

\paragraph{语料来源与预处理}
研究对象为人民银行按季发布的《中国货币政策执行报告》，样本覆盖 2001 年第一季度至
2026 年第二季度共 102 期，全部取自人民银行官网公开页面。每期报告以 PDF 形式存档，
平均 54 页（最短 11 页、最长 73 页）。文本抽取采用 pypdf，随后做三类清洗：去除页眉
页脚与单独成行的页码；合并被换行截断的句子（2001--2003 年的 PDF 在中文字符之间插入了
空格，需先删除）；剔除附表与专栏目录之后的统计附录。

清洗后单期正文平均 25481 个汉字，中位数 26551，最短 9762（2003Q2），最长 35208。
早期报告篇幅明显短于近年，这一点在解释分段长度时需要留意。

\paragraph{情感词典}
情绪度量采用姜富伟等（2020）发布的中文金融情感词典。该词典针对中文财经文本构建，
经人工筛选与词向量扩充，包含积极词 3338 个、消极词 5890 个，使用需引用
Jiang et al.（2019）与姜富伟等（2020）。考虑到本文采用子串匹配识别情感词，
单字词容易在无关词语中被误命中（例如“好”出现在“好处”“良好”之外的搭配中），
因此剔除 4 个单字积极词与 17 个单字消极词；去重后实际使用积极词 3190 个、
消极词 5604 个。

\paragraph{情感单元与计分}
计分以句子为基本单元。把段落按句号、问号、感叹号与分号切分，在每个句子内部自左向右
扫描，优先匹配较长的词（最多四字），命中后跳过该词长度以避免重复计数。若命中词之前
三个字符内出现否定词（“不”“未”“无”“没”“非”等），则翻转其极性。
单期情绪定义为
\begin{equation}
\text{tone}=\frac{N_{\text{pos}}-N_{\text{neg}}}{N_{\text{pos}}+N_{\text{neg}}},
\end{equation}
其中 $N_{\text{pos}}$、$N_{\text{neg}}$ 分别为翻转极性后的积极词与消极词计数。
该比值落在 $[-1,1]$，取值越高表示措辞越偏宽松。

\paragraph{政策指引段与宏观回顾段的切分}
既有研究指出，报告中对市场起作用的主要是面向未来的政策指引，而非对已公布宏观数据的
回顾（姜富伟等，2021）。为区分二者，本文按报告的章节结构切分：以“货币政策趋势”
（2005 年以后）、“预测与展望”或“预测和展望”（2002Q4--2004Q1）、“下一阶段……”
（2001--2002 年部分报告）作为锚点，锚点之后的正文记为政策指引段（记为 P 段），
之前的部分记为宏观回顾段（记为 R 段）。锚点取全文最后一次出现的位置，以避免命中目录；
若 P 段之后出现“附录”“附表”等标题，则在其处截断。

分段结果见表~\ref{tab:corpus}。P 段平均 3216 字，约占全文的 12.7\%，与报告的章节
安排一致。R 段与 P 段的情绪相关系数为 0.375，说明二者相关但并不同质，分段是必要的。

\begin{table}[htbp]
\centering
\small
\begin{threeparttable}
\caption{报告语料与分段统计}
\label{tab:corpus}
\begin{tabular}{lcccccc}
\toprule
 & 均值 & 标准差 & 最小值 & 中位数 & 最大值 & 观测数 \\
\midrule
全文字数            & 25481 & 4833 & 9762  & 26551 & 35208 & 102 \\
宏观回顾段字数（R） & 22264 & 4436 & 1045  & 22858 & 29765 & 102 \\
政策指引段字数（P） & 3216  & 1737 & 0     & 3122  & 15300 & 102 \\
P 段占全文比重      & 0.127 & 0.089 & 0.000 & 0.120 & 0.936 & 102 \\
P 段情绪（tone\_P） & 0.632 & 0.087 & 0.392 & 0.634 & 0.853 & 99  \\
R 段情绪（tone\_R） & 0.524 & 0.085 & 0.275 & 0.530 & 0.739 & 102 \\
\bottomrule
\end{tabular}
\begin{tablenotes}[flushleft]
\footnotesize
\item 字数指清洗后的汉字数。tone\_P 与 tone\_R 为未标准化的原始比值。
R 段与 P 段情绪的相关系数为 0.375。
\end{tablenotes}
\end{threeparttable}
\end{table}

\paragraph{三项质量说明}
第一，词典覆盖率。政策指引段每千字平均命中 103.8 个情感词，宏观回顾段为 65.0 个，
说明词典对政策表述的识别能力并不弱于对经济描述。

第二，分段失败的处理。102 期中有 3 期（2003Q2、2003Q3、2004Q1）因报告结构差异未能
定位到可靠锚点，未生成 P 段情绪，不纳入回归。另有 1 期（2003Q1）分段后 P 段占比达
93.6\%，明显异常，复核系文本排版造成锚点错位。该期情绪值处于样本中部，剔除后主要
结论不变（见第 5 章稳健性检验），故在主回归中保留并予以说明。

第三，与既有研究的关系。本文的分段依据报告章节，而姜富伟等（2021）按句子内容分类，
两者并不完全等价。这是本文与原文在度量上的主要差异，也是本文结果在方向上与之不完全
一致的可能原因之一，第 6 章将对此作进一步讨论。

\subsection{行业收益与市场组合}

\paragraph{行情数据}
行业行情来自 akshare 提供的申万一级行业指数日线，共 37 个代码，覆盖
1999 年 12 月 30 日至 2026 年 11 月 11 日，合计 6498 个交易日。数据字段包括开高低收
与成交量额。

需要说明的是，37 个代码并非同一套分类。申万于 2014 年与 2021 年两次修订行业分类，
新代码的历史被重置，因此各行业的可用起点分为三档（表~\ref{tab:industry}）。
档 A 的 22 个行业中，5 个（801060、801090、801100、801190、801220）在 2014 年 2 月
停止更新，1 个（801020）在 2021 年 12 月终止，其余 16 个贯穿全样本。

\begin{table}[htbp]
\centering
\small
\begin{threeparttable}
\caption{申万一级行业的数据起点}
\label{tab:industry}
\begin{tabular}{lcl}
\toprule
起点 & 行业数 & 说明 \\
\midrule
1999-12-30 & 22 & 其中 5 个于 2014 年 2 月停止，1 个于 2021 年 12 月终止 \\
2014-02-21 & 12 & 申万 2014 版分类新增或更名的行业 \\
2021-12-13 & 3  & 石油石化、环保、美容护理 \\
\midrule
连续覆盖全样本 & 16 & 本文构造市场组合所使用者 \\
2014 年后覆盖 & 28 & 稳健性检验所使用者 \\
\bottomrule
\end{tabular}
\begin{tablenotes}[flushleft]
\footnotesize
\item 行业分类以申万一级为准，代码与名称的对应关系见附录。
\end{tablenotes}
\end{threeparttable}
\end{table}

\paragraph{市场组合的构造}
本文以 16 个连续行业的等权组合作为市场代理：对每个交易日，取当日有行情的行业日收益
的平均值，再累乘为指数。选择等权而非市值加权，主要出于两点考虑：一是收盘价数据中
不含流通市值，无法可靠加权；二是等权组合与后文的行业面板口径一致，便于对照。

这一处理有代价。等权组合偏向中小市值行业，对大盘蓝筹的代表性弱于沪深 300。
本文未能取得覆盖 2001 年以来的沪深 300 日线，因而以行业组合替代，并在 5.5 节用
28 个行业的口径检验结论对行业集合的敏感性。读者在解读效应量级时应考虑这一差异。

\paragraph{收益与波动的定义}
设市场指数在交易日 $t$ 的取值为 $I_t$。发布后一日收益为 $R^{(1)}=I_{t_1}/I_{t_0}-1$，
五日收益为 $R^{(5)}=I_{t_5}/I_{t_0}-1$，其中 $t_0$、$t_1$、$t_5$ 的定义见 3.3 节。
累计异常收益取 $I_{t_1}/I_{t_{-1}}-1$，即包含发布前一个交易日在内的三日窗口。

波动指标有两个。一是事件日绝对收益 $|R^{(0)}|=|I_{t_0}/I_{t_{-1}}-1|$，计算简便但
无法区分单边移动与日内震荡。二是日内振幅，定义为各行业当日最高价与最低价之差除以前
收盘价，再对 16 个行业取平均。相比绝对收益，振幅不受跳空缺口影响，是更干净的波动代理。
此外，本文以发布前 60 个交易日市场收益的标准差衡量事前波动，用于控制波动聚集。

\subsection{事件窗口与时间对齐}

报告的发布日与报告期并不一致。102 期报告全部在报告期结束后的下一个季度发布，
例如 2026 年第一季度报告于 2026 年 5 月 11 日发布。若把报告期收益与当期情绪放在
同一行，等于用 5 月已知的信息解释 1 至 3 月的行情，构成前视偏差。因此本文以发布日
为事件起点，全部收益与波动指标都在发布之后度量。

事件日 $t_0$ 取发布日当日；若发布日不是交易日，则顺延至其后第一个交易日。
102 期中有 6 期需要顺延：2010Q4、2013Q4、2015Q4、2019Q3、2020Q1、2022Q1。
其中 2015Q4 最为特殊，报告于 2016 年 2 月 6 日（周六）发布，其后是春节休市，
第一个交易日为 2016 年 2 月 15 日，距发布已 9 天。报告发布与市场开盘之间的信息
积累会集中释放，该期的事件窗收益不宜与其他期直接比较，本文在稳健性检验中单独标注。

以 $t_0$ 为起点，$t_1$、$t_5$、$t_{-1}$ 分别为其后第 1 个、第 5 个交易日与之前第
1 个交易日。样本中有 10 期的事件窗跨越 9 个以上自然日（春节或国庆），$t_0$ 至 $t_5$
的自然日跨度中位数为 7 天。所有窗口均按交易日推进，不使用自然日，以免在长假期间
产生空窗口。

\subsection{变量定义与样本}

表~\ref{tab:var} 汇总主要变量。情绪变量均作标准化处理，均值为 0、标准差为 1，
因此回归系数可直接读作情绪每变动一个标准差的效应。制度断点变量 $\text{post}_{2014}$
在 2014 年及以后取 1，交互项为标准化情绪与该虚拟变量的乘积。

\begin{table}[htbp]
\centering
\small
\begin{threeparttable}
\caption{主要变量定义}
\label{tab:var}
\begin{tabular}{lll}
\toprule
变量 & 定义 & 说明 \\
\midrule
tone\_P        & 政策指引段情绪，标准化                 & 分母为零时记为缺失 \\
tone\_R        & 宏观回顾段情绪，标准化                 & 同上 \\
$R^{(1)}$      & 发布后一日市场收益                     & $I_{t_1}/I_{t_0}-1$ \\
$R^{(5)}$      & 发布后五日市场收益                     & $I_{t_5}/I_{t_0}-1$ \\
CAR            & 发布前后三日累计收益                   & $I_{t_1}/I_{t_{-1}}-1$ \\
$|R^{(0)}|$    & 事件日绝对收益                         & 波动代理之一 \\
Range          & 事件日日内振幅                         & 16 行业 $(H-L)/C_{-1}$ 的均值 \\
PreVol         & 事前 60 日波动                         & 发布前市场收益标准差 \\
$\text{post}_{2014}$ & 2014 年及以后为 1               & 制度断点虚拟变量 \\
\bottomrule
\end{tabular}
\begin{tablenotes}[flushleft]
\footnotesize
\item 收益与振幅均为小数，回归中系数按基点（$10^{-4}$）报告。
\end{tablenotes}
\end{threeparttable}
\end{table}

经上述处理，可用于回归的样本为 99 期，覆盖 2001Q1 至 2026Q2；其中 98 期可计算
日内振幅（1 期因行情缺失）。2014 年及以后的子样本为 50 期。

\paragraph{数据与代码可得性}
报告原文取自人民银行官网公开页面，行业行情取自 akshare 公开接口，均为公开可得数据，
不涉及脱敏问题。情感词典的使用遵循其说明页的引用要求。全部清洗、度量与估计脚本随论文
一并提供，按说明可复现正文表格。

% =====================================================================
% 第 4 章  研究设计
% ---------------------------------------------------------------------
% 编译说明同前：ctexart 下用 \section，ctexbook/report 下改为 \chapter。
% 依赖宏包：booktabs、threeparttable、amsmath
% =====================================================================

\section{研究设计}

本章说明识别策略。核心问题是：报告级情绪每期只有一个取值，如何在行业面板中识别它的
影响；以及在这一约束下，本文选择什么样的估计方式、做了哪些检验来支撑结论。

\subsection{识别上的一个约束}

设 $i$ 表示行业、$t$ 表示季度，考虑常见的双向固定效应设定
\begin{equation}
y_{i,t}=\alpha+\beta\,\text{tone}_t+\gamma'X_{i,t}+\mu_i+\lambda_t+\varepsilon_{i,t},
\label{eq:twoway}
\end{equation}
其中 $\mu_i$ 为行业固定效应，$\lambda_t$ 为季度固定效应。问题在于
$\text{tone}_t$ 在同一季度内对所有行业取同一个值，它是 $t$ 的函数；而
$\lambda_t$ 已经为每个季度提供了一个独立的截距。二者完全共线，$\beta$ 无法识别，
估计软件会直接把 $\text{tone}_t$ 标记为共线并剔除。

这不是数据量不足或样本期太短造成的，而是模型设定与变量层级的必然结果：任何只在时间
维度上变化的解释变量，在加入时间固定效应后都不可识别。行业面板中确实存在大量
$i\times t$ 层面的变异，但情绪没有这一层面的变异。

由此产生两条可行的路径。第一条是放弃时间固定效应，直接在时间序列上识别，代价是无法
控制同期的共同冲击；第二条是保留固定效应，转而估计情绪与行业特征的交互项
$\text{tone}_t\times Z_i$，此时交互项在 $i\times t$ 层面有变异，斜率可以识别。本文采用
第一条路径，理由在下一小节说明；第二条路径作为后续扩展在 4.5 节讨论。

\subsection{以事件窗口作为识别来源}

本文把识别建立在报告发布日附近的市场反应上，而不是季度的平均状态。这样做的考虑有三点。

其一，报告的影响在时间上是集中的。政策信息在发布当天进入价格，若用整个季度的收益去
回归，等于把事件效应与季度内的其他信息混在一起，稀释了信号。事件窗把观测限制在信息
到达的附近，噪音相对可控。

其二，事件窗内的比较是“同一时点、同一批行业”，共同冲击被时间维度自动吸收。所有行业
在同一交易日面对同一份报告，行业间的差异不需要靠固定效应剔除，样本本身已经是同日
对照。

其三，事件设计使“报告期”与“发布期”的错位得到处理。第 3 章已经说明，报告总是在报告期
结束后的下一季度发布；以发布日为事件起点，可以避免用事后信息解释事前行情。

需要明确的是，这一设计识别的是发布日前后的条件关联，而非严格的因果效应。报告内容与
当期经济状况相关，措辞的变化可能同时反映基本面，因此系数的解释应限定为“给定发布时点，
情绪较高与随后的市场反应之间的关系”。本文在 4.5 节讨论这一限制及其缓解办法。

\subsection{基准模型}

对每一期报告 $e$，估计
\begin{equation}
Y_e=\alpha+\beta\,\text{tone}_{P,e}+\gamma'X_e+\varepsilon_e,
\label{eq:baseline}
\end{equation}
其中 $Y_e$ 为事件窗内的市场收益或波动指标，取一日收益 $R^{(1)}$、五日收益 $R^{(5)}$、
三日累计收益 CAR、事件日绝对收益 $|R^{(0)}|$ 与日内振幅 Range 五种；
$\text{tone}_{P,e}$ 为标准化后的政策指引段情绪；$X_e$ 为控制变量。

控制变量的选择针对一个具体的混淆。情绪较高的报告往往出现在市场相对平静的时期，
若不加控制，情绪与低波动之间的相关可能只是“平静期写得更乐观”的反映。为此在波动
方程中加入事前 60 个交易日的市场收益标准差 $\text{PreVol}_e$，用以吸收波动聚集。
在双段回归中同时纳入宏观回顾段情绪 $\text{tone}_{R,e}$，用于区分“政策指引”与
“整篇语气”两种解释。

标准误采用 Newey--West 估计，滞后阶数取 4。事件间隔约一个季度，事件窗之间不重叠，
因此主要考虑的是残差的自相关与异方差，而非重叠窗口带来的相关性。

与既有研究常用的 EGARCH 设定相比，本文选择事件窗回归出于两点考虑。一是事件窗直接
对应信息到达的时点，条件方差模型则把波动视为状态的连续演化，难以把效应归到某一期
报告上。二是本文需要检验断点，事件窗回归可以通过交互项直接完成，而 GARCH 族模型的
参数时变性需要另作设定。为弥补这一差异，本文同时使用绝对收益与日内振幅两个波动代理，
并在稳健性中报告五日窗口的结果。

\subsection{制度断点的设定}

第 5 章关心的第二个问题是效应是否稳定。设定的出发点是样本期内货币政策框架的调整：
2013 年的流动性紧张、2014 年利率市场化的推进、2015 年汇率形成机制的调整，以及报告
自身章节结构在 2015 年前后的改写。这些变化都可能改变市场解读文本的方式。

据此在式~\eqref{eq:baseline} 中加入断点虚拟变量及其交互项：
\begin{equation}
Y_e=\alpha+\beta_1\,\text{tone}_{P,e}
+\beta_2\,\text{tone}_{P,e}\times\text{post}_{2014}
+\beta_3\,\text{post}_{2014}+\varepsilon_e,
\label{eq:break}
\end{equation}
其中 $\text{post}_{2014}$ 在 2014 年及以后取 1。$\beta_1$ 为断点前的斜率，
$\beta_1+\beta_2$ 为断点后的斜率，$\beta_2$ 即效应的变化量。

断点年份由制度节点确定，而非在数据中搜索得到。为说明这一点，本文把断点依次替换为
2008、2010、2012、2014、2016、2018 年重新估计，报告交互项系数与 $t$ 值的变化。
若结论只在某一年份成立，则说明断点是人为挑选的；若多个相邻年份方向一致，则说明变化
是渐进的。第 5 章表 5.4 报告了这一结果。

\subsection{稳健性设计与识别假设}

\paragraph{四类稳健性检验}
第一，更换市场口径。把等权组合由 16 个行业扩展为 28 个行业，检验结论是否依赖行业集合。
第二，更换情绪词典。以自建货币政策领域词典重新度量情绪，观察系数方向与显著性如何变化。
第三，影响点检验。逐一剔除 2015 年（股市异常波动）、2020 年（疫情冲击）与 2003Q1
（分段异常），并对收益作 1\%/99\% 缩尾，检验结论是否由少数观测支撑。
第四，置换检验。将情绪与事件随机配对 500 次，统计真实 $t$ 值在随机分布中的位置，
这一检验不依赖残差分布的正态假设。

\paragraph{统计功效}
在解读“不显著”之前，需要知道这条设计能识别多大的效应。以一日收益方程为例，系数标准误
约为 25 个基点；在 5\% 显著性水平、80\% 功效下，可识别的最小效应约为 70 个基点，
而日频市场收益本身的标准差接近 190 个基点。这意味着小于 70 个基点的效应在本样本中
无法被可靠识别，不显著的系数不能直接解读为“效应为零”。本文在报告收益方程结果时
一并说明这一限制。

\paragraph{尚不能排除的问题}
其一，同期公告的干扰。货币政策操作、金融统计数据、物价数据等常在相近时点发布，
事件窗内可能同时包含其他信息。本文未能逐日控制这些公告，只能通过窗口长度与安慰剂
检验间接缓解。

其二，发布时点不明。元数据中不含发布的具体时刻，无法判断报告是在盘中还是收盘后
发布。本文以发布当日为 $t_0$；由于无法判断确切时刻，第 5 章同时报告三日累计收益与
五日窗口的结果，用以检验结论对窗口起点的敏感性。

其三，情绪度量的误差。词典法只能识别词面极性，无法处理反讽、条件句与语境依赖；
测量误差会使系数向零收缩，因而本文报告的效应量级更可能是下界。第 5 章以替代词典
检验说明这一偏误的方向与幅度。

其四，反向因果。报告措辞可能与当期经济状况共同决定，本文的断点与置换检验只能缓解
而不能排除这一可能。因此全文的表述避免使用“导致”，而采用“与……相关”“显著降低”
这类描述性措辞。

\paragraph{关于行业层面}
本文的市场组合由行业等权构造，28 行业口径用作稳健性，但并未单独估计“情绪对哪些行业
影响更大”。原因在于 4.1 节所述的识别约束：在季度面板中，情绪与时间固定效应共线，
即使加入行业特征交互项，也需要行业特征本身在样本期内有足够变异，而现有数据中可用的
行业特征（估值、杠杆、利率敏感度）覆盖期偏短，难以支撑稳定估计。这一扩展留待后续
研究，本文对此不作结论。

% =====================================================================
% 第 5 章  实证结果
% ---------------------------------------------------------------------
% 编译说明：
%   主文档建议用 ctexart（\documentclass[12pt]{ctexart}），本文件作为
%   % =====================================================================
% 第 5 章  实证结果
% ---------------------------------------------------------------------
% 编译说明：
%   主文档建议用 ctexart（\documentclass[12pt]{ctexart}），本文件作为
%   % =====================================================================
% 第 5 章  实证结果
% ---------------------------------------------------------------------
% 编译说明：
%   主文档建议用 ctexart（\documentclass[12pt]{ctexart}），本文件作为
%   \input{ch5_empirical} 引入；若主文档是 ctexbook/report，把下面的
%   \section{实证结果} 改成 \chapter{实证结果}，并把 \subsection 降一级。
%   依赖宏包：booktabs、threeparttable、amsmath、graphicx
%   表格编号按"第 5 章"手工写定，若章节顺序调整需同步修改。
% =====================================================================

\section{实证结果}

本章报告文本情绪对市场收益与波动的影响。基准样本为 2001 年第一季度至
2026 年第二季度共 102 期《货币政策执行报告》，其中 99 期可完整切分出政策指引段
与宏观回顾段的情绪指标。市场收益由 16 个连续可得的申万一级行业等权组合构造，
并以 28 个行业口径作为稳健性对照。情绪变量均经标准化处理，回归系数因此可解释为
情绪每变动一个标准差的效应；标准误采用 Newey--West 异方差自相关稳健估计，
滞后阶数取 4。

\subsection{描述性统计}

表~\ref{tab:desc} 列出主要变量的分布。政策指引段情绪（tone\_P）与宏观回顾段情绪
（tone\_R）的相关系数为 0.375，二者相关但并非同一信息的重复度量，这为后续的
分段检验提供了必要的前提。事件日市场收益的均值为 19 个基点，标准差 186 个基点；
日内振幅均值为 2.09\%，意味着报告发布当日最高价与最低价之差平均约为前收盘价的
2\%。事前波动（发布前 60 个交易日市场收益的标准差）均值为 1.46\%。

\begin{table}[htbp]
\centering
\small
\begin{threeparttable}
\caption{描述性统计}
\label{tab:desc}
\begin{tabular}{lcccccc}
\toprule
变量 & 均值 & 标准差 & 最小值 & 中位数 & 最大值 & 观测数 \\
\midrule
政策指引段情绪（tone\_P，标准化） & 0.000 & 1.005 & $-2.773$ & 0.029 & 2.558 & 99 \\
宏观回顾段情绪（tone\_R，标准化） & 0.000 & 1.005 & $-2.928$ & 0.068 & 2.535 & 102 \\
发布后一日收益（\%）              & 0.19  & 1.86  & $-7.60$ & 0.30  & 4.63  & 102 \\
发布后五日收益（\%）              & 0.52  & 3.59  & $-12.59$ & 0.71 & 11.58 & 102 \\
事件日日内振幅（\%）              & 2.09  & 0.92  & 0.80  & 1.88  & 4.76  & 101 \\
事件日绝对收益（\%）              & 0.97  & 0.92  & 0.00  & 0.66  & 3.70  & 102 \\
事前 60 日波动（\%）              & 1.46  & 0.64  & 0.61  & 1.33  & 3.51  & 102 \\
\bottomrule
\end{tabular}
\begin{tablenotes}[flushleft]
\footnotesize
\item 样本为 2001Q1--2026Q2 共 102 期报告；收益与振幅由 16 个申万一级行业的等权组合计算。
\end{tablenotes}
\end{threeparttable}
\end{table}

值得先说明的是情绪与市场波动状态的原始相关性。tone\_P 与事件日日内振幅的相关系数为
$-0.302$，与事件日绝对收益为 $-0.254$，与事前波动的相关系数为 $-0.232$。最后一个
数字提示，情绪较高的报告往往出现在市场相对平静的时期，若不加以控制，很容易把
``情绪降低波动''与``情绪在高波动期偏低''混为一谈。因此在回归中同时报告含事前波动
控制的设定，并以此作为主要判断依据。

\subsection{政策指引情绪与市场波动}

表~\ref{tab:vol} 是本章的主结果。第 (1) 列以事件日绝对收益衡量波动，tone\_P 的系数为
$-23.4$ 个基点（$t=-2.53$）；加入事前波动后（第 (2) 列），系数降至 $-16.6$，显著性
减弱至 10\% 水平，说明有一部分关联确实来自波动状态的持续性。由于绝对收益无法区分
``价格单边移动''与``日内大幅震荡''，第 (3)、(4) 列改用日内振幅，系数分别为 $-27.8$
（$t=-2.58$）与 $-18.5$（$t=-2.43$），控制事前波动后依然在 5\% 水平显著。

第 (5) 列同时纳入回顾段情绪。tone\_R 的系数为 $-2.1$（$t=-0.20$），几乎为零，而
tone\_P 维持在 $-17.8$（$t=-2.27$）。这一对比有明确含义：显著影响波动的并不是
``整篇报告听起来是否乐观''，而是其中面向未来的政策指引部分。第 (6) 列把振幅扩展到
发布后五个交易日，系数为 $-12.1$（$t=-2.03$），效应并未随窗口延长而累积，更接近于
一次性的信息释放而非持续影响。

就经济意义而言，以振幅均值 2.09\% 为参照，$-18.5$ 个基点相当于振幅下降约 9\%；
以标准差衡量，一个标准差的情绪变动对应约 0.2 个标准差的振幅变化。这一量级不算大，
但在 99 个事件、日频噪声显著的条件下已属不易识别。

需要提醒的是，这一幅度对是否包含危机年份较为敏感。剔除 2007--2010 年后，系数降至
$-15.5$（$t=-2.07$），约为基准的一半（见 5.5 节）。方向与显著性不变，但把 $-28$
读作一个稳定不变的效应量并不恰当，更稳妥的表述是：效应量在 15 至 28 个基点之间，
取决于样本是否包含金融危机时期。

分组视角补充了一个细节。把情绪按三分位分组后，低情绪组的平均振幅为 2.53\%，中、高
两组分别为 1.75\% 与 1.99\%（图 5.2）。若把样本切成五组，均值依次为
2.83\%、2.07\%、1.65\%、2.13\%、1.78\%，同样不是单调下降。也就是说，效应主要来自
情绪偏低的一端：措辞异常谨慎的时期波动显著更高，而中高情绪之间并无梯度。这与
``沟通收敛分歧''的解释相容——市场在意的未必是措辞有多乐观，而是它是否透露出不安。

\begin{table}[htbp]
\centering
\small
\begin{threeparttable}
\caption{政策指引段情绪与市场波动}
\label{tab:vol}
\begin{tabular}{lcccccc}
\toprule
 & (1) & (2) & (3) & (4) & (5) & (6) \\
 & $|r|$ & $|r|$ & 振幅 & 振幅 & 振幅 & 五日振幅 \\
\midrule
tone\_P        & $-23.4^{**}$ & $-16.6^{*}$ & $-27.8^{***}$ & $-18.5^{**}$ & $-17.8^{**}$ & $-12.1^{**}$ \\
                & $(-2.53)$ & $(-1.96)$ & $(-2.58)$ & $(-2.43)$ & $(-2.27)$ & $(-2.03)$ \\
tone\_R        & & & & & $-2.1$ & \\
                & & & & & $(-0.20)$ & \\
事前波动        & & 46.0$^{***}$ & & 61.9$^{***}$ & 61.1$^{***}$ & 80.0$^{***}$ \\
                & & $(4.10)$ & & $(4.24)$ & $(4.14)$ & $(4.16)$ \\
\midrule
观测数          & 99 & 99 & 98 & 98 & 98 & 99 \\
$R^2$           & 0.065 & 0.162 & 0.091 & 0.267 & 0.268 & 0.355 \\
\bottomrule
\end{tabular}
\begin{tablenotes}[flushleft]
\footnotesize
\item 被解释变量为事件日市场波动，单位为百分比；情绪与事前波动的系数均已换算为基点
（1 个基点 $=0.01$ 个百分点），其中事前波动的系数读作``事前波动每上升 1 个百分点''。
括号内为 Newey--West 稳健 $t$ 值。$^{*}$、$^{**}$、$^{***}$ 分别表示 10\%、5\%、1\% 显著性水平。
\end{tablenotes}
\end{threeparttable}
\end{table}

\subsection{制度断点：2014 年前后}

样本期内货币政策框架发生了明显变化：2013 年流动性紧张、2014 年利率市场化加速、
2015 年汇率形成机制调整，报告自身的章节结构也在 2015 年前后改写。这些变化提示
情绪的作用方式未必稳定。表~\ref{tab:break} 在方程中引入 2014 年虚拟变量及其与
情绪的交互项，直接检验这一可能。

收益方程中，交互项系数为 $-75.6$ 个基点（$t=-2.11$），五日窗口扩大到 $-181.4$
（$t=-2.65$），即 2014 年以前情绪对收益方向无显著影响（点估计 $+31.6$，$t=0.97$），
2014 年以后转为负向。波动方程则呈现相反的变化：2014 年以前的系数为 $-35.0$
（$t=-3.10$），交互项为 $+26.6$（$t=1.59$），说明波动渠道在后期减弱，但未完全消失。

需要说明的是断点的选择并非由数据挑选而来，而是对应上述制度节点。为检验这一处理是否
稳健，表~\ref{tab:breaksens} 将断点依次替换为 2008、2010、2012、2014、2016、2018 年。
2008 年作断点时交互项接近于零，其余各年份的系数方向一致，其中 2010、2012、2014、2018
年显著，2016 年偏弱。这说明变化是渐进的，把它归因于某一个特定年份并不合适，但
``2010 年代之后效应改变''这一判断是站得住的。

\begin{table}[htbp]
\centering
\small
\begin{threeparttable}
\caption{制度断点检验}
\label{tab:break}
\begin{tabular}{lcccc}
\toprule
 & (1) & (2) & (3) & (4) \\
 & 一日收益 & 五日收益 & 日内振幅 & $|r|$ \\
\midrule
tone\_P              & $+31.6$   & $+73.7$   & $-35.0^{***}$ & $-32.7^{***}$ \\
                      & $(0.97)$  & $(1.22)$  & $(-3.10)$     & $(-2.97)$ \\
tone\_P $\times$ 2014后 & $-75.6^{**}$ & $-181.4^{***}$ & $+26.6$ & $+25.1$ \\
                      & $(-2.11)$ & $(-2.65)$ & $(1.59)$      & $(1.53)$ \\
\midrule
观测数                & 99 & 99 & 98 & 99 \\
$R^2$                 & 0.076 & 0.077 & 0.190 & 0.084 \\
\bottomrule
\end{tabular}
\begin{tablenotes}[flushleft]
\footnotesize
\item 各列均控制 2014 年虚拟变量本身；系数单位为基点，括号内为 Newey--West 稳健 $t$ 值。
\end{tablenotes}
\end{threeparttable}
\end{table}

\begin{table}[htbp]
\centering
\small
\begin{threeparttable}
\caption{断点年份敏感性（交互项系数与 $t$ 值）}
\label{tab:breaksens}
\begin{tabular}{lccc}
\toprule
断点年份 & 一日收益 & 五日收益 & 日内振幅 \\
\midrule
2008 & $-5.9\ (-0.14)$   & $-1.6\ (-0.02)$    & $+1.8\ (0.10)$ \\
2010 & $-86.6^{**}\ (-2.41)$ & $-193.1^{***}\ (-2.69)$ & $+34.5^{**}\ (2.04)$ \\
2012 & $-76.6^{**}\ (-2.20)$ & $-183.5^{***}\ (-2.74)$ & $+27.4^{*}\ (1.71)$ \\
2014 & $-75.6^{**}\ (-2.11)$ & $-181.4^{***}\ (-2.65)$ & $+26.6\ (1.59)$ \\
2016 & $-57.7\ (-1.47)$  & $-134.9^{*}\ (-1.88)$   & $+34.4^{**}\ (2.07)$ \\
2018 & $-64.4^{*}\ (-1.82)$ & $-164.9^{**}\ (-2.29)$ & $+16.2\ (0.79)$ \\
\bottomrule
\end{tabular}
\begin{tablenotes}[flushleft]
\footnotesize
\item 每行对应一次单独的回归，仅改变断点年份，其余设定与表~\ref{tab:break} 相同。
\end{tablenotes}
\end{threeparttable}
\end{table}

\subsection{2014 年之后的方向性效应}

既然断点检验显示后期情绪与收益转为负相关，这里单独报告 2014 年以后的子样本。
表~\ref{tab:post14} 显示，tone\_P 对发布后一日收益的系数为 $-44.0$ 个基点
（$t=-2.89$），五日窗口为 $-107.8$（$t=-3.30$）；同时，波动类方程在子样本中
全部不显著。两组结果互为印证：后期的定价体现在方向上，而不在波动上。

对这一结果需要保持克制。首先，负号的含义并非``情绪越宽松市场越差''这么简单，
更可能反映宽松措辞集中出现在基本面走弱的时点，市场将政策宽松解读为经济前景的
负面信号。其次，该结果对危机年份高度敏感：剔除 2015 年后，一日窗口的交互项由
$-75.6$（$t=-2.11$）降至 $-59.0$（$t=-1.62$）；剔除 2008--2009 年后降至 $-45.5$
（$t=-1.56$）。任一段被剔除，5\% 水平的显著性即不复存在。因此本文不把它作为独立
结论，仅记录符号的方向性变化。

\begin{table}[htbp]
\centering
\small
\begin{threeparttable}
\caption{2014 年以后的子样本}
\label{tab:post14}
\begin{tabular}{lcccc}
\toprule
 & (1) & (2) & (3) & (4) \\
 & 一日收益 & 五日收益 & 一日·双段 & 日内振幅 \\
\midrule
tone\_P & $-44.0^{***}$ & $-107.8^{***}$ & $-47.2^{**}$ & $-8.4$ \\
        & $(-2.89)$     & $(-3.30)$      & $(-2.36)$    & $(-0.68)$ \\
tone\_R &               &                & $+5.9$       & \\
        &               &                & $(0.28)$     & \\
\midrule
观测数  & 50 & 50 & 50 & 49 \\
\bottomrule
\end{tabular}
\begin{tablenotes}[flushleft]
\footnotesize
\item 子样本为 2014Q1--2026Q2。系数单位为基点，括号内为 Newey--West 稳健 $t$ 值。
\end{tablenotes}
\end{threeparttable}
\end{table}

\subsection{稳健性检验}

表~\ref{tab:robust} 汇总四个方面的检验，表~\ref{tab:robust2} 补充三项针对性检验。

\emph{市场口径}。把等权组合由 16 个行业扩展为 28 个行业，主结果基本不变：
日内振幅系数 $-26.7$（$t=-2.58$），断点交互项 $-73.1$（$t=-2.01$），五日交互项
$-183.6$（$t=-2.71$）。结论不依赖于行业集合的选择。

\emph{词典}。以本文自建的货币政策领域词典（积极 159 词、消极 139 词）重新度量情绪，
各系数方向与主结果一致，但量级缩小、显著性下降：绝对收益方程仍显著（$-19.1$，
$t=-2.43$），日内振幅方程降至 $-7.8$（$t=-0.66$），断点交互项降至 $-60.5$
（$t=-1.25$）。该词典每千字仅命中约 68 个情感词，主词典为 104 个，覆盖率差异恰好
对应系数的等比例收缩，符合测量误差导致衰减的预期。因此这一结果是``方向稳健、精度不足''，
而非``结论被推翻''。

\emph{影响点}。逐一剔除 2015 年（股市异常波动）与 2020 年（疫情冲击），并对收益做
1\%/99\% 缩尾。波动类结果对上述处理均不敏感，日内振幅系数稳定在 $-25.8$ 至 $-27.9$
之间。断点交互项则较为敏感：剔除 2015 年后，一日窗口的交互项由 $-75.6$（$t=-2.11$）
降至 $-59.0$（$t=-1.62$），显著性消失；五日窗口降至 $-152.5$（$t=-2.29$），仍在 5\%
水平显著。剔除 2020 年与收益缩尾后，一日交互项分别为 $-73.0$ 与 $-74.4$，保持在 5\%
水平。可见后期方向效应确有相当一部分由 2015 年的数期报告贡献，这解释了上文为何
采取保留态度，也说明该效应在更长窗口上比在单日窗口上更稳定。

\emph{置换检验}。将情绪序列与事件随机配对 500 次，观察真实 $t$ 值在随机分布中的位置。
日内振幅方程的置换 $p$ 值为 0.020，加入回顾段与事前波动后为 0.068，断点交互项为 0.058。
置换检验不依赖分布假设，其结果与稳健 $t$ 检验基本一致：波动结果在 5\% 水平成立，
断点结果处于边际。

\emph{危机期敏感性}。2007--2009 年的措辞与波动同时处于极端，是本文最担心的混淆来源。
剔除 2007--2010 年后，振幅方程系数降至 $-15.5$（$t=-2.07$），剔除 2008--2009 年后为
$-18.8$（$t=-2.34$），加入年份趋势后为 $-20.7$（$t=-2.18$）。方向与显著性保持不变，
幅度约为基准的一半。断点交互项则不同：剔除 2008--2009 年后降至 $-45.5$（$t=-1.56$），
不再显著。可见波动结果稳健但幅度受危机年份影响，方向结果则更依赖少数极端时期。

\emph{安慰剂窗口}。把事件日整体平移 $\pm$5、$\pm$10、$\pm$20 个交易日重估。平移 $\pm$10
天以上时，振幅方程的系数不再显著（$t$ 值均小于 1.1），断点交互项也不再成立，说明效应
与发布时点相关。但平移 $-5$ 天时系数仍有 $-26.6$（$t=-1.79$），保留了约九成的幅度。
这表明效应并非严格的单日冲击，而是集中在发布前后约一周的窗口内，与报告部分内容提前
释放的解释相容。

\emph{政策操作控制}。一个自然的担心是，文本情绪只是货币政策操作的代理。本文整理了三类
总量操作——全面降准 24 次、1 年期 LPR 下调 12 次、7 天逆回购利率下调 12 次——按报告期
汇总为季度操作强度，加入振幅方程。情绪系数由 $-27.8$ 变为 $-27.3$（$t=-2.53$），加入
事前波动后为 $-16.7$（$t=-2.17$）；三类操作本身均不显著，与情绪的相关性也只在 $-0.01$
到 $-0.21$ 之间。断点交互项由 $-75.5$ 变为 $-73.9$（$t=-2.12$）。这说明波动效应并非
由降准降息驱动。

\begin{table}[htbp]
\centering
\small
\begin{threeparttable}
\caption{稳健性检验汇总}
\label{tab:robust}
\begin{tabular}{lcccc}
\toprule
 & 主结果 & 28 行业 & 替代词典 & 剔除 2015 \\
\midrule
\multicolumn{5}{l}{\emph{面板 A：日内振幅（系数，$t$ 值）}} \\
tone\_P & $-27.8^{***}$ & $-26.7^{***}$ & $-7.8$ & $-25.8^{**}$ \\
        & $(-2.58)$ & $(-2.58)$ & $(-0.66)$ & $(-2.19)$ \\
\midrule
\multicolumn{5}{l}{\emph{面板 B：断点交互项（系数，$t$ 值）}} \\
tone\_P $\times$ 2014后（一日） & $-75.6^{**}$ & $-73.1^{**}$ & $-60.5$ & $-59.0$ \\
        & $(-2.11)$ & $(-2.01)$ & $(-1.25)$ & $(-1.62)$ \\
tone\_P $\times$ 2014后（五日） & $-181.4^{***}$ & $-183.6^{***}$ & $-109.3$ & $-152.5^{**}$ \\
        & $(-2.65)$ & $(-2.71)$ & $(-1.25)$ & $(-2.29)$ \\
\midrule
\multicolumn{5}{l}{\emph{面板 C：置换检验 $p$ 值}} \\
 & 0.020 & — & — & — \\
\bottomrule
\end{tabular}
\begin{tablenotes}[flushleft]
\footnotesize
\item ``替代词典''指以自建货币政策领域词典重新计算情绪指标的结果；``剔除 2015''指去掉 2015 年
四期报告后重估。面板 C 仅报告日内振幅方程的置换 $p$ 值。
\end{tablenotes}
\end{threeparttable}
\end{table}

\begin{table}[htbp]
\centering
\small
\begin{threeparttable}
\caption{补充检验：危机期、安慰剂窗口与政策操作}
\label{tab:robust2}
\begin{tabular}{lcc}
\toprule
设定 & 日内振幅（系数，$t$ 值） & 断点交互项（系数，$t$ 值） \\
\midrule
基准 & $-27.8^{***}$ $(-2.58)$ & $-75.5^{**}$ $(-2.11)$ \\
剔除 2007--2010 & $-15.5^{**}$ $(-2.07)$ & $-36.6$ $(-1.27)$ \\
剔除 2008--2009 & $-18.8^{**}$ $(-2.34)$ & $-45.5$ $(-1.56)$ \\
加入年份趋势 & $-20.7^{**}$ $(-2.18)$ & — \\
事件日前移 20 个交易日 & $+8.5$ $(0.46)$ & $+9.6$ $(0.40)$ \\
事件日后移 20 个交易日 & $-12.5$ $(-1.21)$ & $-35.1$ $(-1.51)$ \\
控制三类政策操作 & $-27.3^{**}$ $(-2.53)$ & $-73.9^{**}$ $(-2.12)$ \\
\bottomrule
\end{tabular}
\begin{tablenotes}[flushleft]
\footnotesize
\item 系数单位为基点，括号内为 Newey--West 稳健 $t$ 值。安慰剂检验按平移后的日期
重新计算振幅与收益，情绪变量不变。政策操作指全面降准、1 年期 LPR 与 7 天逆回购利率，
按报告期汇总为季度强度。``—''表示该项未单独估计。
\end{tablenotes}
\end{threeparttable}
\end{table}

综合来看，本章的证据支持一个有限但明确的判断：政策指引段的文本情绪确实影响市场，
但其作用首先体现在波动而非收益方向上；方向性效应在 2014 年之后才出现，且稳健性
弱于波动结果。这一区分对理解央行沟通的作用机制是重要的——若只报告全样本的平均效应，
两个方向相反的结果会相互抵消，从而得出``沟通无效''的错误结论。
 引入；若主文档是 ctexbook/report，把下面的
%   \section{实证结果} 改成 \chapter{实证结果}，并把 \subsection 降一级。
%   依赖宏包：booktabs、threeparttable、amsmath、graphicx
%   表格编号按"第 5 章"手工写定，若章节顺序调整需同步修改。
% =====================================================================

\section{实证结果}

本章报告文本情绪对市场收益与波动的影响。基准样本为 2001 年第一季度至
2026 年第二季度共 102 期《货币政策执行报告》，其中 99 期可完整切分出政策指引段
与宏观回顾段的情绪指标。市场收益由 16 个连续可得的申万一级行业等权组合构造，
并以 28 个行业口径作为稳健性对照。情绪变量均经标准化处理，回归系数因此可解释为
情绪每变动一个标准差的效应；标准误采用 Newey--West 异方差自相关稳健估计，
滞后阶数取 4。

\subsection{描述性统计}

表~\ref{tab:desc} 列出主要变量的分布。政策指引段情绪（tone\_P）与宏观回顾段情绪
（tone\_R）的相关系数为 0.375，二者相关但并非同一信息的重复度量，这为后续的
分段检验提供了必要的前提。事件日市场收益的均值为 19 个基点，标准差 186 个基点；
日内振幅均值为 2.09\%，意味着报告发布当日最高价与最低价之差平均约为前收盘价的
2\%。事前波动（发布前 60 个交易日市场收益的标准差）均值为 1.46\%。

\begin{table}[htbp]
\centering
\small
\begin{threeparttable}
\caption{描述性统计}
\label{tab:desc}
\begin{tabular}{lcccccc}
\toprule
变量 & 均值 & 标准差 & 最小值 & 中位数 & 最大值 & 观测数 \\
\midrule
政策指引段情绪（tone\_P，标准化） & 0.000 & 1.005 & $-2.773$ & 0.029 & 2.558 & 99 \\
宏观回顾段情绪（tone\_R，标准化） & 0.000 & 1.005 & $-2.928$ & 0.068 & 2.535 & 102 \\
发布后一日收益（\%）              & 0.19  & 1.86  & $-7.60$ & 0.30  & 4.63  & 102 \\
发布后五日收益（\%）              & 0.52  & 3.59  & $-12.59$ & 0.71 & 11.58 & 102 \\
事件日日内振幅（\%）              & 2.09  & 0.92  & 0.80  & 1.88  & 4.76  & 101 \\
事件日绝对收益（\%）              & 0.97  & 0.92  & 0.00  & 0.66  & 3.70  & 102 \\
事前 60 日波动（\%）              & 1.46  & 0.64  & 0.61  & 1.33  & 3.51  & 102 \\
\bottomrule
\end{tabular}
\begin{tablenotes}[flushleft]
\footnotesize
\item 样本为 2001Q1--2026Q2 共 102 期报告；收益与振幅由 16 个申万一级行业的等权组合计算。
\end{tablenotes}
\end{threeparttable}
\end{table}

值得先说明的是情绪与市场波动状态的原始相关性。tone\_P 与事件日日内振幅的相关系数为
$-0.302$，与事件日绝对收益为 $-0.254$，与事前波动的相关系数为 $-0.232$。最后一个
数字提示，情绪较高的报告往往出现在市场相对平静的时期，若不加以控制，很容易把
``情绪降低波动''与``情绪在高波动期偏低''混为一谈。因此在回归中同时报告含事前波动
控制的设定，并以此作为主要判断依据。

\subsection{政策指引情绪与市场波动}

表~\ref{tab:vol} 是本章的主结果。第 (1) 列以事件日绝对收益衡量波动，tone\_P 的系数为
$-23.4$ 个基点（$t=-2.53$）；加入事前波动后（第 (2) 列），系数降至 $-16.6$，显著性
减弱至 10\% 水平，说明有一部分关联确实来自波动状态的持续性。由于绝对收益无法区分
``价格单边移动''与``日内大幅震荡''，第 (3)、(4) 列改用日内振幅，系数分别为 $-27.8$
（$t=-2.58$）与 $-18.5$（$t=-2.43$），控制事前波动后依然在 5\% 水平显著。

第 (5) 列同时纳入回顾段情绪。tone\_R 的系数为 $-2.1$（$t=-0.20$），几乎为零，而
tone\_P 维持在 $-17.8$（$t=-2.27$）。这一对比有明确含义：显著影响波动的并不是
``整篇报告听起来是否乐观''，而是其中面向未来的政策指引部分。第 (6) 列把振幅扩展到
发布后五个交易日，系数为 $-12.1$（$t=-2.03$），效应并未随窗口延长而累积，更接近于
一次性的信息释放而非持续影响。

就经济意义而言，以振幅均值 2.09\% 为参照，$-18.5$ 个基点相当于振幅下降约 9\%；
以标准差衡量，一个标准差的情绪变动对应约 0.2 个标准差的振幅变化。这一量级不算大，
但在 99 个事件、日频噪声显著的条件下已属不易识别。

需要提醒的是，这一幅度对是否包含危机年份较为敏感。剔除 2007--2010 年后，系数降至
$-15.5$（$t=-2.07$），约为基准的一半（见 5.5 节）。方向与显著性不变，但把 $-28$
读作一个稳定不变的效应量并不恰当，更稳妥的表述是：效应量在 15 至 28 个基点之间，
取决于样本是否包含金融危机时期。

分组视角补充了一个细节。把情绪按三分位分组后，低情绪组的平均振幅为 2.53\%，中、高
两组分别为 1.75\% 与 1.99\%（图 5.2）。若把样本切成五组，均值依次为
2.83\%、2.07\%、1.65\%、2.13\%、1.78\%，同样不是单调下降。也就是说，效应主要来自
情绪偏低的一端：措辞异常谨慎的时期波动显著更高，而中高情绪之间并无梯度。这与
``沟通收敛分歧''的解释相容——市场在意的未必是措辞有多乐观，而是它是否透露出不安。

\begin{table}[htbp]
\centering
\small
\begin{threeparttable}
\caption{政策指引段情绪与市场波动}
\label{tab:vol}
\begin{tabular}{lcccccc}
\toprule
 & (1) & (2) & (3) & (4) & (5) & (6) \\
 & $|r|$ & $|r|$ & 振幅 & 振幅 & 振幅 & 五日振幅 \\
\midrule
tone\_P        & $-23.4^{**}$ & $-16.6^{*}$ & $-27.8^{***}$ & $-18.5^{**}$ & $-17.8^{**}$ & $-12.1^{**}$ \\
                & $(-2.53)$ & $(-1.96)$ & $(-2.58)$ & $(-2.43)$ & $(-2.27)$ & $(-2.03)$ \\
tone\_R        & & & & & $-2.1$ & \\
                & & & & & $(-0.20)$ & \\
事前波动        & & 46.0$^{***}$ & & 61.9$^{***}$ & 61.1$^{***}$ & 80.0$^{***}$ \\
                & & $(4.10)$ & & $(4.24)$ & $(4.14)$ & $(4.16)$ \\
\midrule
观测数          & 99 & 99 & 98 & 98 & 98 & 99 \\
$R^2$           & 0.065 & 0.162 & 0.091 & 0.267 & 0.268 & 0.355 \\
\bottomrule
\end{tabular}
\begin{tablenotes}[flushleft]
\footnotesize
\item 被解释变量为事件日市场波动，单位为百分比；情绪与事前波动的系数均已换算为基点
（1 个基点 $=0.01$ 个百分点），其中事前波动的系数读作``事前波动每上升 1 个百分点''。
括号内为 Newey--West 稳健 $t$ 值。$^{*}$、$^{**}$、$^{***}$ 分别表示 10\%、5\%、1\% 显著性水平。
\end{tablenotes}
\end{threeparttable}
\end{table}

\subsection{制度断点：2014 年前后}

样本期内货币政策框架发生了明显变化：2013 年流动性紧张、2014 年利率市场化加速、
2015 年汇率形成机制调整，报告自身的章节结构也在 2015 年前后改写。这些变化提示
情绪的作用方式未必稳定。表~\ref{tab:break} 在方程中引入 2014 年虚拟变量及其与
情绪的交互项，直接检验这一可能。

收益方程中，交互项系数为 $-75.6$ 个基点（$t=-2.11$），五日窗口扩大到 $-181.4$
（$t=-2.65$），即 2014 年以前情绪对收益方向无显著影响（点估计 $+31.6$，$t=0.97$），
2014 年以后转为负向。波动方程则呈现相反的变化：2014 年以前的系数为 $-35.0$
（$t=-3.10$），交互项为 $+26.6$（$t=1.59$），说明波动渠道在后期减弱，但未完全消失。

需要说明的是断点的选择并非由数据挑选而来，而是对应上述制度节点。为检验这一处理是否
稳健，表~\ref{tab:breaksens} 将断点依次替换为 2008、2010、2012、2014、2016、2018 年。
2008 年作断点时交互项接近于零，其余各年份的系数方向一致，其中 2010、2012、2014、2018
年显著，2016 年偏弱。这说明变化是渐进的，把它归因于某一个特定年份并不合适，但
``2010 年代之后效应改变''这一判断是站得住的。

\begin{table}[htbp]
\centering
\small
\begin{threeparttable}
\caption{制度断点检验}
\label{tab:break}
\begin{tabular}{lcccc}
\toprule
 & (1) & (2) & (3) & (4) \\
 & 一日收益 & 五日收益 & 日内振幅 & $|r|$ \\
\midrule
tone\_P              & $+31.6$   & $+73.7$   & $-35.0^{***}$ & $-32.7^{***}$ \\
                      & $(0.97)$  & $(1.22)$  & $(-3.10)$     & $(-2.97)$ \\
tone\_P $\times$ 2014后 & $-75.6^{**}$ & $-181.4^{***}$ & $+26.6$ & $+25.1$ \\
                      & $(-2.11)$ & $(-2.65)$ & $(1.59)$      & $(1.53)$ \\
\midrule
观测数                & 99 & 99 & 98 & 99 \\
$R^2$                 & 0.076 & 0.077 & 0.190 & 0.084 \\
\bottomrule
\end{tabular}
\begin{tablenotes}[flushleft]
\footnotesize
\item 各列均控制 2014 年虚拟变量本身；系数单位为基点，括号内为 Newey--West 稳健 $t$ 值。
\end{tablenotes}
\end{threeparttable}
\end{table}

\begin{table}[htbp]
\centering
\small
\begin{threeparttable}
\caption{断点年份敏感性（交互项系数与 $t$ 值）}
\label{tab:breaksens}
\begin{tabular}{lccc}
\toprule
断点年份 & 一日收益 & 五日收益 & 日内振幅 \\
\midrule
2008 & $-5.9\ (-0.14)$   & $-1.6\ (-0.02)$    & $+1.8\ (0.10)$ \\
2010 & $-86.6^{**}\ (-2.41)$ & $-193.1^{***}\ (-2.69)$ & $+34.5^{**}\ (2.04)$ \\
2012 & $-76.6^{**}\ (-2.20)$ & $-183.5^{***}\ (-2.74)$ & $+27.4^{*}\ (1.71)$ \\
2014 & $-75.6^{**}\ (-2.11)$ & $-181.4^{***}\ (-2.65)$ & $+26.6\ (1.59)$ \\
2016 & $-57.7\ (-1.47)$  & $-134.9^{*}\ (-1.88)$   & $+34.4^{**}\ (2.07)$ \\
2018 & $-64.4^{*}\ (-1.82)$ & $-164.9^{**}\ (-2.29)$ & $+16.2\ (0.79)$ \\
\bottomrule
\end{tabular}
\begin{tablenotes}[flushleft]
\footnotesize
\item 每行对应一次单独的回归，仅改变断点年份，其余设定与表~\ref{tab:break} 相同。
\end{tablenotes}
\end{threeparttable}
\end{table}

\subsection{2014 年之后的方向性效应}

既然断点检验显示后期情绪与收益转为负相关，这里单独报告 2014 年以后的子样本。
表~\ref{tab:post14} 显示，tone\_P 对发布后一日收益的系数为 $-44.0$ 个基点
（$t=-2.89$），五日窗口为 $-107.8$（$t=-3.30$）；同时，波动类方程在子样本中
全部不显著。两组结果互为印证：后期的定价体现在方向上，而不在波动上。

对这一结果需要保持克制。首先，负号的含义并非``情绪越宽松市场越差''这么简单，
更可能反映宽松措辞集中出现在基本面走弱的时点，市场将政策宽松解读为经济前景的
负面信号。其次，该结果对危机年份高度敏感：剔除 2015 年后，一日窗口的交互项由
$-75.6$（$t=-2.11$）降至 $-59.0$（$t=-1.62$）；剔除 2008--2009 年后降至 $-45.5$
（$t=-1.56$）。任一段被剔除，5\% 水平的显著性即不复存在。因此本文不把它作为独立
结论，仅记录符号的方向性变化。

\begin{table}[htbp]
\centering
\small
\begin{threeparttable}
\caption{2014 年以后的子样本}
\label{tab:post14}
\begin{tabular}{lcccc}
\toprule
 & (1) & (2) & (3) & (4) \\
 & 一日收益 & 五日收益 & 一日·双段 & 日内振幅 \\
\midrule
tone\_P & $-44.0^{***}$ & $-107.8^{***}$ & $-47.2^{**}$ & $-8.4$ \\
        & $(-2.89)$     & $(-3.30)$      & $(-2.36)$    & $(-0.68)$ \\
tone\_R &               &                & $+5.9$       & \\
        &               &                & $(0.28)$     & \\
\midrule
观测数  & 50 & 50 & 50 & 49 \\
\bottomrule
\end{tabular}
\begin{tablenotes}[flushleft]
\footnotesize
\item 子样本为 2014Q1--2026Q2。系数单位为基点，括号内为 Newey--West 稳健 $t$ 值。
\end{tablenotes}
\end{threeparttable}
\end{table}

\subsection{稳健性检验}

表~\ref{tab:robust} 汇总四个方面的检验，表~\ref{tab:robust2} 补充三项针对性检验。

\emph{市场口径}。把等权组合由 16 个行业扩展为 28 个行业，主结果基本不变：
日内振幅系数 $-26.7$（$t=-2.58$），断点交互项 $-73.1$（$t=-2.01$），五日交互项
$-183.6$（$t=-2.71$）。结论不依赖于行业集合的选择。

\emph{词典}。以本文自建的货币政策领域词典（积极 159 词、消极 139 词）重新度量情绪，
各系数方向与主结果一致，但量级缩小、显著性下降：绝对收益方程仍显著（$-19.1$，
$t=-2.43$），日内振幅方程降至 $-7.8$（$t=-0.66$），断点交互项降至 $-60.5$
（$t=-1.25$）。该词典每千字仅命中约 68 个情感词，主词典为 104 个，覆盖率差异恰好
对应系数的等比例收缩，符合测量误差导致衰减的预期。因此这一结果是``方向稳健、精度不足''，
而非``结论被推翻''。

\emph{影响点}。逐一剔除 2015 年（股市异常波动）与 2020 年（疫情冲击），并对收益做
1\%/99\% 缩尾。波动类结果对上述处理均不敏感，日内振幅系数稳定在 $-25.8$ 至 $-27.9$
之间。断点交互项则较为敏感：剔除 2015 年后，一日窗口的交互项由 $-75.6$（$t=-2.11$）
降至 $-59.0$（$t=-1.62$），显著性消失；五日窗口降至 $-152.5$（$t=-2.29$），仍在 5\%
水平显著。剔除 2020 年与收益缩尾后，一日交互项分别为 $-73.0$ 与 $-74.4$，保持在 5\%
水平。可见后期方向效应确有相当一部分由 2015 年的数期报告贡献，这解释了上文为何
采取保留态度，也说明该效应在更长窗口上比在单日窗口上更稳定。

\emph{置换检验}。将情绪序列与事件随机配对 500 次，观察真实 $t$ 值在随机分布中的位置。
日内振幅方程的置换 $p$ 值为 0.020，加入回顾段与事前波动后为 0.068，断点交互项为 0.058。
置换检验不依赖分布假设，其结果与稳健 $t$ 检验基本一致：波动结果在 5\% 水平成立，
断点结果处于边际。

\emph{危机期敏感性}。2007--2009 年的措辞与波动同时处于极端，是本文最担心的混淆来源。
剔除 2007--2010 年后，振幅方程系数降至 $-15.5$（$t=-2.07$），剔除 2008--2009 年后为
$-18.8$（$t=-2.34$），加入年份趋势后为 $-20.7$（$t=-2.18$）。方向与显著性保持不变，
幅度约为基准的一半。断点交互项则不同：剔除 2008--2009 年后降至 $-45.5$（$t=-1.56$），
不再显著。可见波动结果稳健但幅度受危机年份影响，方向结果则更依赖少数极端时期。

\emph{安慰剂窗口}。把事件日整体平移 $\pm$5、$\pm$10、$\pm$20 个交易日重估。平移 $\pm$10
天以上时，振幅方程的系数不再显著（$t$ 值均小于 1.1），断点交互项也不再成立，说明效应
与发布时点相关。但平移 $-5$ 天时系数仍有 $-26.6$（$t=-1.79$），保留了约九成的幅度。
这表明效应并非严格的单日冲击，而是集中在发布前后约一周的窗口内，与报告部分内容提前
释放的解释相容。

\emph{政策操作控制}。一个自然的担心是，文本情绪只是货币政策操作的代理。本文整理了三类
总量操作——全面降准 24 次、1 年期 LPR 下调 12 次、7 天逆回购利率下调 12 次——按报告期
汇总为季度操作强度，加入振幅方程。情绪系数由 $-27.8$ 变为 $-27.3$（$t=-2.53$），加入
事前波动后为 $-16.7$（$t=-2.17$）；三类操作本身均不显著，与情绪的相关性也只在 $-0.01$
到 $-0.21$ 之间。断点交互项由 $-75.5$ 变为 $-73.9$（$t=-2.12$）。这说明波动效应并非
由降准降息驱动。

\begin{table}[htbp]
\centering
\small
\begin{threeparttable}
\caption{稳健性检验汇总}
\label{tab:robust}
\begin{tabular}{lcccc}
\toprule
 & 主结果 & 28 行业 & 替代词典 & 剔除 2015 \\
\midrule
\multicolumn{5}{l}{\emph{面板 A：日内振幅（系数，$t$ 值）}} \\
tone\_P & $-27.8^{***}$ & $-26.7^{***}$ & $-7.8$ & $-25.8^{**}$ \\
        & $(-2.58)$ & $(-2.58)$ & $(-0.66)$ & $(-2.19)$ \\
\midrule
\multicolumn{5}{l}{\emph{面板 B：断点交互项（系数，$t$ 值）}} \\
tone\_P $\times$ 2014后（一日） & $-75.6^{**}$ & $-73.1^{**}$ & $-60.5$ & $-59.0$ \\
        & $(-2.11)$ & $(-2.01)$ & $(-1.25)$ & $(-1.62)$ \\
tone\_P $\times$ 2014后（五日） & $-181.4^{***}$ & $-183.6^{***}$ & $-109.3$ & $-152.5^{**}$ \\
        & $(-2.65)$ & $(-2.71)$ & $(-1.25)$ & $(-2.29)$ \\
\midrule
\multicolumn{5}{l}{\emph{面板 C：置换检验 $p$ 值}} \\
 & 0.020 & — & — & — \\
\bottomrule
\end{tabular}
\begin{tablenotes}[flushleft]
\footnotesize
\item ``替代词典''指以自建货币政策领域词典重新计算情绪指标的结果；``剔除 2015''指去掉 2015 年
四期报告后重估。面板 C 仅报告日内振幅方程的置换 $p$ 值。
\end{tablenotes}
\end{threeparttable}
\end{table}

\begin{table}[htbp]
\centering
\small
\begin{threeparttable}
\caption{补充检验：危机期、安慰剂窗口与政策操作}
\label{tab:robust2}
\begin{tabular}{lcc}
\toprule
设定 & 日内振幅（系数，$t$ 值） & 断点交互项（系数，$t$ 值） \\
\midrule
基准 & $-27.8^{***}$ $(-2.58)$ & $-75.5^{**}$ $(-2.11)$ \\
剔除 2007--2010 & $-15.5^{**}$ $(-2.07)$ & $-36.6$ $(-1.27)$ \\
剔除 2008--2009 & $-18.8^{**}$ $(-2.34)$ & $-45.5$ $(-1.56)$ \\
加入年份趋势 & $-20.7^{**}$ $(-2.18)$ & — \\
事件日前移 20 个交易日 & $+8.5$ $(0.46)$ & $+9.6$ $(0.40)$ \\
事件日后移 20 个交易日 & $-12.5$ $(-1.21)$ & $-35.1$ $(-1.51)$ \\
控制三类政策操作 & $-27.3^{**}$ $(-2.53)$ & $-73.9^{**}$ $(-2.12)$ \\
\bottomrule
\end{tabular}
\begin{tablenotes}[flushleft]
\footnotesize
\item 系数单位为基点，括号内为 Newey--West 稳健 $t$ 值。安慰剂检验按平移后的日期
重新计算振幅与收益，情绪变量不变。政策操作指全面降准、1 年期 LPR 与 7 天逆回购利率，
按报告期汇总为季度强度。``—''表示该项未单独估计。
\end{tablenotes}
\end{threeparttable}
\end{table}

综合来看，本章的证据支持一个有限但明确的判断：政策指引段的文本情绪确实影响市场，
但其作用首先体现在波动而非收益方向上；方向性效应在 2014 年之后才出现，且稳健性
弱于波动结果。这一区分对理解央行沟通的作用机制是重要的——若只报告全样本的平均效应，
两个方向相反的结果会相互抵消，从而得出``沟通无效''的错误结论。
 引入；若主文档是 ctexbook/report，把下面的
%   \section{实证结果} 改成 \chapter{实证结果}，并把 \subsection 降一级。
%   依赖宏包：booktabs、threeparttable、amsmath、graphicx
%   表格编号按"第 5 章"手工写定，若章节顺序调整需同步修改。
% =====================================================================

\section{实证结果}

本章报告文本情绪对市场收益与波动的影响。基准样本为 2001 年第一季度至
2026 年第二季度共 102 期《货币政策执行报告》，其中 99 期可完整切分出政策指引段
与宏观回顾段的情绪指标。市场收益由 16 个连续可得的申万一级行业等权组合构造，
并以 28 个行业口径作为稳健性对照。情绪变量均经标准化处理，回归系数因此可解释为
情绪每变动一个标准差的效应；标准误采用 Newey--West 异方差自相关稳健估计，
滞后阶数取 4。

\subsection{描述性统计}

表~\ref{tab:desc} 列出主要变量的分布。政策指引段情绪（tone\_P）与宏观回顾段情绪
（tone\_R）的相关系数为 0.375，二者相关但并非同一信息的重复度量，这为后续的
分段检验提供了必要的前提。事件日市场收益的均值为 19 个基点，标准差 186 个基点；
日内振幅均值为 2.09\%，意味着报告发布当日最高价与最低价之差平均约为前收盘价的
2\%。事前波动（发布前 60 个交易日市场收益的标准差）均值为 1.46\%。

\begin{table}[htbp]
\centering
\small
\begin{threeparttable}
\caption{描述性统计}
\label{tab:desc}
\begin{tabular}{lcccccc}
\toprule
变量 & 均值 & 标准差 & 最小值 & 中位数 & 最大值 & 观测数 \\
\midrule
政策指引段情绪（tone\_P，标准化） & 0.000 & 1.005 & $-2.773$ & 0.029 & 2.558 & 99 \\
宏观回顾段情绪（tone\_R，标准化） & 0.000 & 1.005 & $-2.928$ & 0.068 & 2.535 & 102 \\
发布后一日收益（\%）              & 0.19  & 1.86  & $-7.60$ & 0.30  & 4.63  & 102 \\
发布后五日收益（\%）              & 0.52  & 3.59  & $-12.59$ & 0.71 & 11.58 & 102 \\
事件日日内振幅（\%）              & 2.09  & 0.92  & 0.80  & 1.88  & 4.76  & 101 \\
事件日绝对收益（\%）              & 0.97  & 0.92  & 0.00  & 0.66  & 3.70  & 102 \\
事前 60 日波动（\%）              & 1.46  & 0.64  & 0.61  & 1.33  & 3.51  & 102 \\
\bottomrule
\end{tabular}
\begin{tablenotes}[flushleft]
\footnotesize
\item 样本为 2001Q1--2026Q2 共 102 期报告；收益与振幅由 16 个申万一级行业的等权组合计算。
\end{tablenotes}
\end{threeparttable}
\end{table}

值得先说明的是情绪与市场波动状态的原始相关性。tone\_P 与事件日日内振幅的相关系数为
$-0.302$，与事件日绝对收益为 $-0.254$，与事前波动的相关系数为 $-0.232$。最后一个
数字提示，情绪较高的报告往往出现在市场相对平静的时期，若不加以控制，很容易把
``情绪降低波动''与``情绪在高波动期偏低''混为一谈。因此在回归中同时报告含事前波动
控制的设定，并以此作为主要判断依据。

\subsection{政策指引情绪与市场波动}

表~\ref{tab:vol} 是本章的主结果。第 (1) 列以事件日绝对收益衡量波动，tone\_P 的系数为
$-23.4$ 个基点（$t=-2.53$）；加入事前波动后（第 (2) 列），系数降至 $-16.6$，显著性
减弱至 10\% 水平，说明有一部分关联确实来自波动状态的持续性。由于绝对收益无法区分
``价格单边移动''与``日内大幅震荡''，第 (3)、(4) 列改用日内振幅，系数分别为 $-27.8$
（$t=-2.58$）与 $-18.5$（$t=-2.43$），控制事前波动后依然在 5\% 水平显著。

第 (5) 列同时纳入回顾段情绪。tone\_R 的系数为 $-2.1$（$t=-0.20$），几乎为零，而
tone\_P 维持在 $-17.8$（$t=-2.27$）。这一对比有明确含义：显著影响波动的并不是
``整篇报告听起来是否乐观''，而是其中面向未来的政策指引部分。第 (6) 列把振幅扩展到
发布后五个交易日，系数为 $-12.1$（$t=-2.03$），效应并未随窗口延长而累积，更接近于
一次性的信息释放而非持续影响。

就经济意义而言，以振幅均值 2.09\% 为参照，$-18.5$ 个基点相当于振幅下降约 9\%；
以标准差衡量，一个标准差的情绪变动对应约 0.2 个标准差的振幅变化。这一量级不算大，
但在 99 个事件、日频噪声显著的条件下已属不易识别。

需要提醒的是，这一幅度对是否包含危机年份较为敏感。剔除 2007--2010 年后，系数降至
$-15.5$（$t=-2.07$），约为基准的一半（见 5.5 节）。方向与显著性不变，但把 $-28$
读作一个稳定不变的效应量并不恰当，更稳妥的表述是：效应量在 15 至 28 个基点之间，
取决于样本是否包含金融危机时期。

分组视角补充了一个细节。把情绪按三分位分组后，低情绪组的平均振幅为 2.53\%，中、高
两组分别为 1.75\% 与 1.99\%（图 5.2）。若把样本切成五组，均值依次为
2.83\%、2.07\%、1.65\%、2.13\%、1.78\%，同样不是单调下降。也就是说，效应主要来自
情绪偏低的一端：措辞异常谨慎的时期波动显著更高，而中高情绪之间并无梯度。这与
``沟通收敛分歧''的解释相容——市场在意的未必是措辞有多乐观，而是它是否透露出不安。

\begin{table}[htbp]
\centering
\small
\begin{threeparttable}
\caption{政策指引段情绪与市场波动}
\label{tab:vol}
\begin{tabular}{lcccccc}
\toprule
 & (1) & (2) & (3) & (4) & (5) & (6) \\
 & $|r|$ & $|r|$ & 振幅 & 振幅 & 振幅 & 五日振幅 \\
\midrule
tone\_P        & $-23.4^{**}$ & $-16.6^{*}$ & $-27.8^{***}$ & $-18.5^{**}$ & $-17.8^{**}$ & $-12.1^{**}$ \\
                & $(-2.53)$ & $(-1.96)$ & $(-2.58)$ & $(-2.43)$ & $(-2.27)$ & $(-2.03)$ \\
tone\_R        & & & & & $-2.1$ & \\
                & & & & & $(-0.20)$ & \\
事前波动        & & 46.0$^{***}$ & & 61.9$^{***}$ & 61.1$^{***}$ & 80.0$^{***}$ \\
                & & $(4.10)$ & & $(4.24)$ & $(4.14)$ & $(4.16)$ \\
\midrule
观测数          & 99 & 99 & 98 & 98 & 98 & 99 \\
$R^2$           & 0.065 & 0.162 & 0.091 & 0.267 & 0.268 & 0.355 \\
\bottomrule
\end{tabular}
\begin{tablenotes}[flushleft]
\footnotesize
\item 被解释变量为事件日市场波动，单位为百分比；情绪与事前波动的系数均已换算为基点
（1 个基点 $=0.01$ 个百分点），其中事前波动的系数读作``事前波动每上升 1 个百分点''。
括号内为 Newey--West 稳健 $t$ 值。$^{*}$、$^{**}$、$^{***}$ 分别表示 10\%、5\%、1\% 显著性水平。
\end{tablenotes}
\end{threeparttable}
\end{table}

\subsection{制度断点：2014 年前后}

样本期内货币政策框架发生了明显变化：2013 年流动性紧张、2014 年利率市场化加速、
2015 年汇率形成机制调整，报告自身的章节结构也在 2015 年前后改写。这些变化提示
情绪的作用方式未必稳定。表~\ref{tab:break} 在方程中引入 2014 年虚拟变量及其与
情绪的交互项，直接检验这一可能。

收益方程中，交互项系数为 $-75.6$ 个基点（$t=-2.11$），五日窗口扩大到 $-181.4$
（$t=-2.65$），即 2014 年以前情绪对收益方向无显著影响（点估计 $+31.6$，$t=0.97$），
2014 年以后转为负向。波动方程则呈现相反的变化：2014 年以前的系数为 $-35.0$
（$t=-3.10$），交互项为 $+26.6$（$t=1.59$），说明波动渠道在后期减弱，但未完全消失。

需要说明的是断点的选择并非由数据挑选而来，而是对应上述制度节点。为检验这一处理是否
稳健，表~\ref{tab:breaksens} 将断点依次替换为 2008、2010、2012、2014、2016、2018 年。
2008 年作断点时交互项接近于零，其余各年份的系数方向一致，其中 2010、2012、2014、2018
年显著，2016 年偏弱。这说明变化是渐进的，把它归因于某一个特定年份并不合适，但
``2010 年代之后效应改变''这一判断是站得住的。

\begin{table}[htbp]
\centering
\small
\begin{threeparttable}
\caption{制度断点检验}
\label{tab:break}
\begin{tabular}{lcccc}
\toprule
 & (1) & (2) & (3) & (4) \\
 & 一日收益 & 五日收益 & 日内振幅 & $|r|$ \\
\midrule
tone\_P              & $+31.6$   & $+73.7$   & $-35.0^{***}$ & $-32.7^{***}$ \\
                      & $(0.97)$  & $(1.22)$  & $(-3.10)$     & $(-2.97)$ \\
tone\_P $\times$ 2014后 & $-75.6^{**}$ & $-181.4^{***}$ & $+26.6$ & $+25.1$ \\
                      & $(-2.11)$ & $(-2.65)$ & $(1.59)$      & $(1.53)$ \\
\midrule
观测数                & 99 & 99 & 98 & 99 \\
$R^2$                 & 0.076 & 0.077 & 0.190 & 0.084 \\
\bottomrule
\end{tabular}
\begin{tablenotes}[flushleft]
\footnotesize
\item 各列均控制 2014 年虚拟变量本身；系数单位为基点，括号内为 Newey--West 稳健 $t$ 值。
\end{tablenotes}
\end{threeparttable}
\end{table}

\begin{table}[htbp]
\centering
\small
\begin{threeparttable}
\caption{断点年份敏感性（交互项系数与 $t$ 值）}
\label{tab:breaksens}
\begin{tabular}{lccc}
\toprule
断点年份 & 一日收益 & 五日收益 & 日内振幅 \\
\midrule
2008 & $-5.9\ (-0.14)$   & $-1.6\ (-0.02)$    & $+1.8\ (0.10)$ \\
2010 & $-86.6^{**}\ (-2.41)$ & $-193.1^{***}\ (-2.69)$ & $+34.5^{**}\ (2.04)$ \\
2012 & $-76.6^{**}\ (-2.20)$ & $-183.5^{***}\ (-2.74)$ & $+27.4^{*}\ (1.71)$ \\
2014 & $-75.6^{**}\ (-2.11)$ & $-181.4^{***}\ (-2.65)$ & $+26.6\ (1.59)$ \\
2016 & $-57.7\ (-1.47)$  & $-134.9^{*}\ (-1.88)$   & $+34.4^{**}\ (2.07)$ \\
2018 & $-64.4^{*}\ (-1.82)$ & $-164.9^{**}\ (-2.29)$ & $+16.2\ (0.79)$ \\
\bottomrule
\end{tabular}
\begin{tablenotes}[flushleft]
\footnotesize
\item 每行对应一次单独的回归，仅改变断点年份，其余设定与表~\ref{tab:break} 相同。
\end{tablenotes}
\end{threeparttable}
\end{table}

\subsection{2014 年之后的方向性效应}

既然断点检验显示后期情绪与收益转为负相关，这里单独报告 2014 年以后的子样本。
表~\ref{tab:post14} 显示，tone\_P 对发布后一日收益的系数为 $-44.0$ 个基点
（$t=-2.89$），五日窗口为 $-107.8$（$t=-3.30$）；同时，波动类方程在子样本中
全部不显著。两组结果互为印证：后期的定价体现在方向上，而不在波动上。

对这一结果需要保持克制。首先，负号的含义并非``情绪越宽松市场越差''这么简单，
更可能反映宽松措辞集中出现在基本面走弱的时点，市场将政策宽松解读为经济前景的
负面信号。其次，该结果对危机年份高度敏感：剔除 2015 年后，一日窗口的交互项由
$-75.6$（$t=-2.11$）降至 $-59.0$（$t=-1.62$）；剔除 2008--2009 年后降至 $-45.5$
（$t=-1.56$）。任一段被剔除，5\% 水平的显著性即不复存在。因此本文不把它作为独立
结论，仅记录符号的方向性变化。

\begin{table}[htbp]
\centering
\small
\begin{threeparttable}
\caption{2014 年以后的子样本}
\label{tab:post14}
\begin{tabular}{lcccc}
\toprule
 & (1) & (2) & (3) & (4) \\
 & 一日收益 & 五日收益 & 一日·双段 & 日内振幅 \\
\midrule
tone\_P & $-44.0^{***}$ & $-107.8^{***}$ & $-47.2^{**}$ & $-8.4$ \\
        & $(-2.89)$     & $(-3.30)$      & $(-2.36)$    & $(-0.68)$ \\
tone\_R &               &                & $+5.9$       & \\
        &               &                & $(0.28)$     & \\
\midrule
观测数  & 50 & 50 & 50 & 49 \\
\bottomrule
\end{tabular}
\begin{tablenotes}[flushleft]
\footnotesize
\item 子样本为 2014Q1--2026Q2。系数单位为基点，括号内为 Newey--West 稳健 $t$ 值。
\end{tablenotes}
\end{threeparttable}
\end{table}

\subsection{稳健性检验}

表~\ref{tab:robust} 汇总四个方面的检验，表~\ref{tab:robust2} 补充三项针对性检验。

\emph{市场口径}。把等权组合由 16 个行业扩展为 28 个行业，主结果基本不变：
日内振幅系数 $-26.7$（$t=-2.58$），断点交互项 $-73.1$（$t=-2.01$），五日交互项
$-183.6$（$t=-2.71$）。结论不依赖于行业集合的选择。

\emph{词典}。以本文自建的货币政策领域词典（积极 159 词、消极 139 词）重新度量情绪，
各系数方向与主结果一致，但量级缩小、显著性下降：绝对收益方程仍显著（$-19.1$，
$t=-2.43$），日内振幅方程降至 $-7.8$（$t=-0.66$），断点交互项降至 $-60.5$
（$t=-1.25$）。该词典每千字仅命中约 68 个情感词，主词典为 104 个，覆盖率差异恰好
对应系数的等比例收缩，符合测量误差导致衰减的预期。因此这一结果是``方向稳健、精度不足''，
而非``结论被推翻''。

\emph{影响点}。逐一剔除 2015 年（股市异常波动）与 2020 年（疫情冲击），并对收益做
1\%/99\% 缩尾。波动类结果对上述处理均不敏感，日内振幅系数稳定在 $-25.8$ 至 $-27.9$
之间。断点交互项则较为敏感：剔除 2015 年后，一日窗口的交互项由 $-75.6$（$t=-2.11$）
降至 $-59.0$（$t=-1.62$），显著性消失；五日窗口降至 $-152.5$（$t=-2.29$），仍在 5\%
水平显著。剔除 2020 年与收益缩尾后，一日交互项分别为 $-73.0$ 与 $-74.4$，保持在 5\%
水平。可见后期方向效应确有相当一部分由 2015 年的数期报告贡献，这解释了上文为何
采取保留态度，也说明该效应在更长窗口上比在单日窗口上更稳定。

\emph{置换检验}。将情绪序列与事件随机配对 500 次，观察真实 $t$ 值在随机分布中的位置。
日内振幅方程的置换 $p$ 值为 0.020，加入回顾段与事前波动后为 0.068，断点交互项为 0.058。
置换检验不依赖分布假设，其结果与稳健 $t$ 检验基本一致：波动结果在 5\% 水平成立，
断点结果处于边际。

\emph{危机期敏感性}。2007--2009 年的措辞与波动同时处于极端，是本文最担心的混淆来源。
剔除 2007--2010 年后，振幅方程系数降至 $-15.5$（$t=-2.07$），剔除 2008--2009 年后为
$-18.8$（$t=-2.34$），加入年份趋势后为 $-20.7$（$t=-2.18$）。方向与显著性保持不变，
幅度约为基准的一半。断点交互项则不同：剔除 2008--2009 年后降至 $-45.5$（$t=-1.56$），
不再显著。可见波动结果稳健但幅度受危机年份影响，方向结果则更依赖少数极端时期。

\emph{安慰剂窗口}。把事件日整体平移 $\pm$5、$\pm$10、$\pm$20 个交易日重估。平移 $\pm$10
天以上时，振幅方程的系数不再显著（$t$ 值均小于 1.1），断点交互项也不再成立，说明效应
与发布时点相关。但平移 $-5$ 天时系数仍有 $-26.6$（$t=-1.79$），保留了约九成的幅度。
这表明效应并非严格的单日冲击，而是集中在发布前后约一周的窗口内，与报告部分内容提前
释放的解释相容。

\emph{政策操作控制}。一个自然的担心是，文本情绪只是货币政策操作的代理。本文整理了三类
总量操作——全面降准 24 次、1 年期 LPR 下调 12 次、7 天逆回购利率下调 12 次——按报告期
汇总为季度操作强度，加入振幅方程。情绪系数由 $-27.8$ 变为 $-27.3$（$t=-2.53$），加入
事前波动后为 $-16.7$（$t=-2.17$）；三类操作本身均不显著，与情绪的相关性也只在 $-0.01$
到 $-0.21$ 之间。断点交互项由 $-75.5$ 变为 $-73.9$（$t=-2.12$）。这说明波动效应并非
由降准降息驱动。

\begin{table}[htbp]
\centering
\small
\begin{threeparttable}
\caption{稳健性检验汇总}
\label{tab:robust}
\begin{tabular}{lcccc}
\toprule
 & 主结果 & 28 行业 & 替代词典 & 剔除 2015 \\
\midrule
\multicolumn{5}{l}{\emph{面板 A：日内振幅（系数，$t$ 值）}} \\
tone\_P & $-27.8^{***}$ & $-26.7^{***}$ & $-7.8$ & $-25.8^{**}$ \\
        & $(-2.58)$ & $(-2.58)$ & $(-0.66)$ & $(-2.19)$ \\
\midrule
\multicolumn{5}{l}{\emph{面板 B：断点交互项（系数，$t$ 值）}} \\
tone\_P $\times$ 2014后（一日） & $-75.6^{**}$ & $-73.1^{**}$ & $-60.5$ & $-59.0$ \\
        & $(-2.11)$ & $(-2.01)$ & $(-1.25)$ & $(-1.62)$ \\
tone\_P $\times$ 2014后（五日） & $-181.4^{***}$ & $-183.6^{***}$ & $-109.3$ & $-152.5^{**}$ \\
        & $(-2.65)$ & $(-2.71)$ & $(-1.25)$ & $(-2.29)$ \\
\midrule
\multicolumn{5}{l}{\emph{面板 C：置换检验 $p$ 值}} \\
 & 0.020 & — & — & — \\
\bottomrule
\end{tabular}
\begin{tablenotes}[flushleft]
\footnotesize
\item ``替代词典''指以自建货币政策领域词典重新计算情绪指标的结果；``剔除 2015''指去掉 2015 年
四期报告后重估。面板 C 仅报告日内振幅方程的置换 $p$ 值。
\end{tablenotes}
\end{threeparttable}
\end{table}

\begin{table}[htbp]
\centering
\small
\begin{threeparttable}
\caption{补充检验：危机期、安慰剂窗口与政策操作}
\label{tab:robust2}
\begin{tabular}{lcc}
\toprule
设定 & 日内振幅（系数，$t$ 值） & 断点交互项（系数，$t$ 值） \\
\midrule
基准 & $-27.8^{***}$ $(-2.58)$ & $-75.5^{**}$ $(-2.11)$ \\
剔除 2007--2010 & $-15.5^{**}$ $(-2.07)$ & $-36.6$ $(-1.27)$ \\
剔除 2008--2009 & $-18.8^{**}$ $(-2.34)$ & $-45.5$ $(-1.56)$ \\
加入年份趋势 & $-20.7^{**}$ $(-2.18)$ & — \\
事件日前移 20 个交易日 & $+8.5$ $(0.46)$ & $+9.6$ $(0.40)$ \\
事件日后移 20 个交易日 & $-12.5$ $(-1.21)$ & $-35.1$ $(-1.51)$ \\
控制三类政策操作 & $-27.3^{**}$ $(-2.53)$ & $-73.9^{**}$ $(-2.12)$ \\
\bottomrule
\end{tabular}
\begin{tablenotes}[flushleft]
\footnotesize
\item 系数单位为基点，括号内为 Newey--West 稳健 $t$ 值。安慰剂检验按平移后的日期
重新计算振幅与收益，情绪变量不变。政策操作指全面降准、1 年期 LPR 与 7 天逆回购利率，
按报告期汇总为季度强度。``—''表示该项未单独估计。
\end{tablenotes}
\end{threeparttable}
\end{table}

综合来看，本章的证据支持一个有限但明确的判断：政策指引段的文本情绪确实影响市场，
但其作用首先体现在波动而非收益方向上；方向性效应在 2014 年之后才出现，且稳健性
弱于波动结果。这一区分对理解央行沟通的作用机制是重要的——若只报告全样本的平均效应，
两个方向相反的结果会相互抵消，从而得出``沟通无效''的错误结论。

% =====================================================================
% 第 6 章  讨论与结论
% ---------------------------------------------------------------------
% 编译说明同前：ctexart 下用 \section，ctexbook/report 下改为 \chapter。
% =====================================================================

\section{讨论与结论}

\subsection{主要发现}

本文用 2001 年至 2026 年的 102 期《货币政策执行报告》检验了文本情绪的定价效应，
得到三点结果。

第一，政策指引段的情绪对发布后的市场收益方向没有稳健影响。全样本中，一日与五日
窗口的系数都接近于零，其标准误对应的可识别区间在数十个基点的量级，说明这一零结果
不是估计失败，而是本样本确实识别不到方向性效应。

第二，同一情绪显著降低了发布日的市场波动。无论以绝对收益还是日内振幅度量，控制
事前波动之后系数仍为负且在 5\% 水平显著，28 个行业的口径给出几乎相同的估计，
置换检验的 $p$ 值为 0.020。

第三，把样本按 2014 年切开后，波动效应集中在前期，方向效应出现在后期且为负。
这一转折对断点年份的选择不敏感，却对剔除 2015 年较为敏感，因此本文把它作为提示性
证据，而非与第二条并列的结论。

\subsection{为什么作用在波动上而不是方向上}

这一结果的直觉并不复杂。政策文件的作用首先是让市场对政策路径形成更一致的判断，
而不是直接告诉市场资产应该值多少钱。当报告传递的指引较为明确时，投资者对未来的
分歧收敛，交易行为随之趋同，价格的单日摆动幅度下降；至于价格往哪个方向走，还取决于
基本面、流动性与估值，单一文件的措辞难以决定。

这一解释与既有证据是相容的。姜富伟等（2021）在四类宽基指数上发现，文本情绪上升
伴随收益上升，文本相似度上升伴随波动下降；两者的共同点在于，文本提供的是关于政策
路径的信息，其对分歧的作用比对水平的作用更直接。王熙和郑梦圆（2021）
把语调指数纳入结构向量自回归，控制信息效应之后“价格之谜”消失，进一步说明央行文本
中确实含有关于未来政策的信息成分，而不只是对已发生事实的复述。

本文的分段结果为这一解释提供了额外支持。若波动下降只是“整篇报告听起来更乐观”的
结果，那么宏观回顾段的情绪应当具有类似作用；实际估计中，回顾段的系数为 $-2.1$
个基点（$t=-0.20$），几乎为零。这提示起作用的是面向未来的部分，与姜富伟等（2021）
关于“政策指引而非宏观回顾”的判断一致。王辉等（2023）把报告拆为前瞻性与回顾性两部分
分别度量，蒋海等（2025）则以机器学习方法提取前瞻性情绪，两项研究的设计都建立在
“前瞻内容承载主要政策信号”这一前提之上。

\subsection{2014 年之后的方向性变化}

方向性效应在 2014 年之后才出现，且符号为负。本文认为有两种互不排斥的解释。

第一种是“坏消息”效应。宽松措辞往往出现在经济下行压力较大的时点，措辞越积极，
越可能意味着政策需要发力的程度越深。市场若把宽松解读为对基本面的确认，报告发布后
反而出现负收益。这一机制与王宇伟等（2025）关于央行言辞沟通非对称性的发现相通：
在他们关于企业家信心的研究中，紧缩表述有效、宽松表述效果有限，说明市场对不同方向
的措辞并非对称处理。

第二种是预期管理的成熟。2014 年之后，央行在报告发布前的沟通渠道明显增多，
政策意图被提前定价，等到报告正式发布，新增信息有限，“利好出尽”式的回调更容易出现。
这一解释可以说明为什么方向效应为负而波动效应同时减弱：当信息已被消化，文件本身
既不能推动价格，也难以收敛分歧。

需要承认的是，本文的数据无法在两种解释之间作出裁决。可用的证据只有一条：剔除 2015
年后方向效应明显减弱（交互项由 $-75.6$ 降至 $-59.0$，$t=-1.62$），剔除 2008--2009 年
后进一步降至 $-45.5$（$t=-1.56$），两者都不再显著。这说明该方向效应相当一部分由金融
危机与股市异常波动期的数期报告贡献，而这类时期恰好是“坏消息”效应最可能显现的时点。
这一点更倾向于第一种解释，但不足以排除第二种。

\subsection{与既有研究的关系}

与姜富伟等（2021）相比，本文的差别有两处。一是样本延长至 2026 年，覆盖了 2020 年
之后的非常规货币政策阶段；二是把被解释变量从收益扩展到波动，并引入制度断点。
结论上，本文确认了“指引段而非回顾段”的判断，但没有复现出显著的方向效应。这一差异
可能来自三方面：分段依据不同（本文按章节，原文按句子内容）、市场代理不同（本文用
行业等权组合，原文用四类宽基指数）、样本期不同。第 3 章已就第一点作了说明。

与周阳、綦岚鑫（2024）相比，两者的差异主要在结果变量。他们构建了专门适配报告的
领域词典，并同时考察股票、债券与外汇市场的收益与波动，发现市场之间存在溢出，
溢出会放大或削弱情绪效应。本文只做股票市场，但在波动维度上得到的结论与之相容：
情绪对二阶矩的作用比对一阶矩更稳定。他们词典中纳入否定词与程度副词的做法，
也为本文后续改进度量提供了方向。

与游庭玮（2026）相比，两者都使用申万一级行业数据，但解释变量不同。他用美联储政策
利率，本文用报告文本；他通过控制市场整体收益识别行业独立冲击，本文则把行业等权组合
直接作为市场代理。两条路线并不冲突，若能取得覆盖期更长的行业特征数据，把文本情绪
与行业暴露度结合，是更自然的下一步。

陈妍玲等（2024）在金融业及其子行业上发现，前瞻性报告的文本情绪显著提高风险溢出，
这与本文“情绪影响二阶矩”的结论方向一致，只是他们关注的是跨市场溢出，本文关注的是
发布日的波动。林建浩和孙乐轩（2025）在综述中提醒，文本度量对模型设定与提示方式较为
敏感，可复现性有待改善。本文以两套词典与置换检验回应了这一提醒，也确实观察到替代
词典下系数衰减的现象。

\subsection{政策含义}

如果沟通的主要作用是收敛分歧，那么它的价值就不应只用“能否拉动指数”来衡量。
在预期管理框架下，一份表述清晰的报告可以降低市场对政策路径的不确定性，减少无谓的
波动，这本身即是政策目标的一部分。本文的估计为这一判断提供了数量上的支持：
情绪每提高一个标准差，发布日的日内振幅下降约 28 个基点，控制事前波动后仍有约
19 个基点。

另一方面，方向性效应的不稳定提示不宜高估单份文件的短期市场影响。2014 年之后方向
效应转为负，且对个别年份敏感，说明市场对报告的解读依赖于当时的宏观环境与预期状态。
对政策制定者而言，这意味着沟通的效果不是恒定的，需要结合所处的政策周期来判断；
对投资者而言，把报告措辞当作方向性交易信号，其可靠性远低于把它当作不确定性指标。

\subsection{局限与后续研究}

本文的局限主要有四。

其一，度量误差。词典法只能识别词面极性，对条件句、反讽与语境依赖的处理能力有限。
替代词典下系数按覆盖率比例收缩，说明误差主要是衰减型的，因而本文报告的效应量级
更可能是下界而非上界。引入带否定词与程度副词的领域词典，或用预训练模型替代词典，
是可行的改进方向。

其二，识别的时间精度。元数据不含发布的具体时刻，事件窗的起点只能以交易日为单位
近似。虽然更换窗口的结果一致，但对盘中反应与隔夜反应的区分仍无法做到。

其三，同期信息的干扰。货币政策操作、金融统计与物价数据常在相近时点公布，事件窗内
可能混入其他信息。本文的处理有两条：一是把降准、LPR 与逆回购操作作为控制变量加入方程，
加入后情绪系数几乎不变；二是把事件日平移若干交易日重做，结果发现效应在 $\pm$10 天以外
消失，但前移 5 天仍有约九成幅度，说明它并非严格的单日冲击。这两点缓解了担忧，但没有
完全消除。

其四，样本与功效。99 个事件、日频数据，能够可靠识别的效应量级在 70 个基点以上，
更细微的效应无法区分。若要把方向效应估计得更精确，需要把事件扩展到其他沟通文本，
例如例会纪要、新闻发布会实录，或改用更高频的数据。

行业层面是本文有意未展开的部分。第 4 章已说明，报告级变量在季度面板中与时间固定
效应共线，要做行业异质性，必须依赖行业特征与情绪的交互，而这要求行业特征在样本期内
有足够变异。现有数据中的估值与财务指标覆盖期偏短，难以支撑稳定估计。若后续能补齐
这部分数据，“哪些行业对政策指引更敏感”仍是一个值得回答的问题。

\subsection{结论}

本文的出发点是一个简单的问题：央行在报告里说的话，市场听进去了吗？答案取决于问的是
哪个方面。就价格方向而言，本文没有找到稳健的证据；就市场分歧而言，证据是清楚的——
政策指引段的措辞越积极，报告发布当日的波动越小，且这一作用并非来自报告的整体语气。

这一区分对理解央行沟通有意义，对评估沟通效果同样有影响。若只估计全样本的平均效应，
2014 年以前的方向性零结果与之后的方向性负效应会相互抵消，得到的结论是“沟通对价格
没有影响”；而真实情况是，沟通的作用先体现在波动上，随时间推移才在方向上有所显现。
把不同时期、不同文本段落视为同质样本，容易掩盖这一结构变化。


\newpage
% =====================================================================
% 参考文献
% ---------------------------------------------------------------------
% 说明：
%   1. 本文正文采用"作者（年份）"的行文方式，未使用 \cite，故此处用
%      thebibliography 直接列出，编译后即为参考文献表。
%   2. 全部条目的出处与卷期已核验，核验过程见 docs/文献引用核验表.md。
%   3. 中文条目按第一作者拼音排序，外文在前。
% =====================================================================

\begin{thebibliography}{99}
\footnotesize

\bibitem{lucca2009}
Lucca D O, Trebbi F. Measuring Central Bank Communication: An Automated
Approach with Application to FOMC Statements. NBER Working Paper No.~15367, 2009.

\bibitem{shapiro2022}
Shapiro A H, Wilson D J. Taking the Fed at Its Word: A New Approach to
Estimating Central Bank Objectives Using Text Analysis.
\textit{Review of Economic Studies}, 2022, 89(5): 2768--2805.

\bibitem{chen2024}
陈妍玲, 唐靖, 郭文伟. 央行沟通与金融系统性风险：促进还是抑制？
\textit{上海立信会计金融学院学报}, 2024, 36(6): 3--20.

\bibitem{ji2026}
纪钧尧, 尤晓彤, 陈旭. 央行货币政策文本情绪对汇率的影响：基于 FinBERT
模型的多国文本情绪挖掘. \textit{世界经济研究}, 2026(7).

\bibitem{jiang2021}
姜富伟, 胡逸驰, 黄楠. 央行货币政策报告文本信息、宏观经济与股票市场.
\textit{金融研究}, 2021(6): 95--113.

\bibitem{jiang2020}
姜富伟, 孟令超, 唐国豪. 媒体文本情绪与股票回报预测.
\textit{经济学（季刊）}, 2020, 已录用待刊.

\bibitem{jiang2019}
Jiang F, Lee J, Martin X, Zhou G. Manager Sentiment and Stock Returns.
\textit{Journal of Financial Economics}, 2019, 132(1): 126--149.

\bibitem{jiang2025}
蒋海, 孙娜, 李欣明. 央行预期管理能缓解银行流动性囤积吗？——基于前瞻性
文本情绪指数的实证检验. \textit{金融研究}, 2025(6): 1--20.

\bibitem{liao2026}
廖海龙. 大模型驱动的情绪因子构建及其在 A 股量化择时中的应用.
江西财经大学硕士学位论文, 2026.

\bibitem{lin2025}
林建浩, 孙乐轩. 大语言模型与经济金融文本分析：基本原理、应用场景与研究展望.
\textit{计量经济学报}, 2025, 5(1).

\bibitem{wang2023}
王辉, 朱家雲, 胡诣聪. 央行预期引导可以降低银行系统性金融风险吗？——基于
市场解读偏离的视角. \textit{金融研究}, 2023(9): 1--19.

\bibitem{wang2021}
王熙, 郑梦圆. 中国货币政策冲击的信息效应与“价格之谜”——基于文本分析的
研究. \textit{山东大学学报（哲学社会科学版）}, 2021(5).

\bibitem{wang2025}
王宇伟, 刘宏雅, 范从来. 央行言辞沟通能否提振企业家信心.
\textit{经济理论与经济管理}, 2025, 45(8): 33--52.

\bibitem{wang2022}
王春芝. 货币政策公告的股价预漂移效应研究——基于风险敏感理论的视角.
南京财经大学硕士学位论文, 2022.

\bibitem{wangyu2025}
汪禹, 吴彬. 情感分析视角下中文政策文本的倾向词典构建研究.
\textit{情报资料工作}, 2025, 46(5).

\bibitem{xu2023}
许魁. 央行货币政策报告文本信息与股票市场——基于分析师报告视角.
山东大学硕士学位论文, 2023.

\bibitem{you2026}
游庭玮. 美联储政策利率对 A 股行业收益的冲击效应研究.
华东政法大学硕士学位论文, 2026.

\bibitem{zhou2024}
周阳, 綦岚鑫. 央行文本情绪具有金融市场效应吗——基于《中国货币政策执行
报告》的研究. \textit{财经科学}, 2024(2).

\bibitem{zhou2018}
周亮. 货币周期对我国股市行业轮动的影响.
\textit{湖北经济学院学报}, 2018(3): 31--37.

\end{thebibliography}


\end{document}
